\documentclass[11pt]{article}

\usepackage[margin=1in]{geometry}
\usepackage{amsmath,amssymb,amsthm}
\usepackage{mathtools}
\usepackage{enumitem}
\usepackage[hidelinks]{hyperref}
\usepackage{microtype}

\newtheorem{theorem}{Theorem}[section]
\newtheorem{lemma}[theorem]{Lemma}
\newtheorem{proposition}[theorem]{Proposition}
\newtheorem{corollary}[theorem]{Corollary}
\newtheorem{conjecture}[theorem]{Conjecture}
\theoremstyle{definition}
\newtheorem{definition}[theorem]{Definition}
\newtheorem{remark}[theorem]{Remark}

\DeclareMathOperator{\lcm}{lcm}
\newcommand{\NN}{\mathbb{N}}
\newcommand{\ZZ}{\mathbb{Z}}
\newcommand{\QQ}{\mathbb{Q}}
\newcommand{\PP}{\mathcal{P}}
\newcommand{\der}{{}'}
\newcommand{\dd}{\,\mathrm{d}}

\title{On the equation $n''=n$ and a problem of Erd\H{o}s and Barbeau \\
       on products of prime--reciprocal sums}
\author{Vico Bonfioli\thanks{\texttt{vico@anvilstack.com}.
  See the acknowledgments for a candid note on the use of AI assistance in this work.}}
\date{\today}

\begin{document}
\maketitle

\begin{abstract}
Erd\H{o}s Problem~\#307 (posed by Barbeau, 1976; recorded in Erd\H{o}s--Graham, 1980) asks
whether there exist two finite sets $P,Q$ of primes with
$\bigl(\sum_{p\in P}\tfrac1p\bigr)\bigl(\sum_{q\in Q}\tfrac1q\bigr)=1$.
Writing $a=\prod P$, $b=\prod Q$, a solution exists if and only if the \emph{arithmetic derivative}
has a two--cycle $a'=b,\ b'=a$ with $a\neq b$; equivalently iff $n''=n$ has a solution that is not a
fixed point $p^{p}$. The two--cycle reformulation is due to Seva (2019) and, independently, Bado
(2026); our starting point is to identify the map with the arithmetic derivative, which ties \#307 to
the Ufnarovski--\AA hlander dynamics: the conjecture that every
$n''=n$--cycle has period one would force \#307 to have answer \textsc{no}, while a \textsc{yes}
would refute it. Exploiting the equivalence we prove a \emph{rigidity lemma} (each prime--reciprocal
sum is automatically in lowest terms, so $N_P=D_Q$, $N_Q=D_P$, and $P\cap Q=\varnothing$) and a
quantitative \emph{barrier}: building on Rosen's bound $|P\cup Q|\ge 59$, any solution has
$\min(\prod P,\prod Q)\ge 2.09\times10^{56}$, improving the best published lower bound for $n''=n$
solutions by about $47$ orders of magnitude. A finite verification over the complete list of
$49{,}961$ admissible $59$--prime supports then closes the $59$--prime level: any solution in fact
has $|P\cup Q|\ge 60$ and $\min(\prod P,\prod Q)>3.50\times10^{57}$. We add a parity dichotomy (a hypothetical all--odd
cycle would need $\ge 1412$ primes and have $\min(\prod P,\prod Q)>10^{2519}$), a proof that longer
cycles are strictly harder ($k=3$ already forces $361{,}139$ primes), and a function--field theorem:
over $\mathbb{F}_q[t]$ the derivative has \emph{no} cycles, so the obstruction over $\mathbb{Z}$ is
purely archimedean. A local congruence anatomy of the cycle equation, the Erd\H{o}s--Wintner fact
that the abundance density $0.0420$ is a genuine constant, and a Th\=abit/Euler--type constructive
calculus (the \emph{Frame
Rule} and a divisor--trick) that faithfully parametrises solutions but, by an exact identity
$\Delta_0=M_0N(1-s_{M_0}s_N)$, provably does \emph{not} reduce the difficulty below \#307 itself
(the natural ``make $\Delta_0$ small'' objective is \#307 in disguise). Recast as a bilinear breeder,
the rule yields the maximal honest consequence --- a conditional reduction $(\mathrm{H})\Rightarrow$
\textsc{yes} to an explicit prime--density hypothesis --- which a quantitative ledger shows is
\emph{not} a feasible finite search. The rigidity result and the barrier are fully formalised in Lean~4 (mathlib), the smallest--primes
extremality being proved rather than assumed; the only non--logical input is a verified evaluation
of the first $59$ primes. The problem remains open; the heuristics weakly favour
existence.
\end{abstract}

\section{Introduction}\label{sec:intro}

Let $P,Q$ be finite, nonempty sets of (distinct) primes. Erd\H{o}s Problem~\#307
\cite{ErdosProblems307}, attributed to Barbeau's 1976 ``Computer challenge corner'' problem
\#477 \cite{Barbeau1976} and listed in Erd\H{o}s--Graham \cite{ErdosGraham1980}, asks:
\begin{quote}
do there exist finite sets of primes $P,Q$ with
$\displaystyle\Bigl(\sum_{p\in P}\frac1p\Bigr)\Bigl(\sum_{q\in Q}\frac1q\Bigr)=1$\,?
\end{quote}
The problem is verifiable in the sense that a single example would settle it. No example is known;
the only published partial fact (a remark of J.~Rosen on MathOverflow \cite{MO320838}, paraphrased
on the problem page) is that any solution needs at least $59$ primes in $P\cup Q$. The known
``solutions'' of Cambie and others solve only the \emph{weakened} version in which $1$ is permitted
as an element, e.g.\ $(1+\tfrac15)(\tfrac12+\tfrac13)=1$; these do not bear on the prime version.

Expanding the product exhibits \#307 as a rigidly shaped Egyptian fraction for $1$:
$1=\sum_{p\in P}\sum_{q\in Q}1/(pq)$, whose denominators are the \emph{rectangular grid} $P\times Q$
of semiprimes --- a \emph{rank--one}, multiplicatively separable representation. The unrestricted
companion --- representing $1$ by reciprocals of distinct two--prime--factor numbers, with no product
structure imposed --- is by contrast \emph{solved}: Johnson (1978) gave a $48$--term decomposition,
Watanabe \cite{Watanabe2020} $47$--term ones, and Graham's survey \cite{GrahamEgyptian} records exactly
this contrast, noting after Barbeau that the product form ``is not known.'' Thus \#307 is precisely the
rank--one restriction of a solved problem. Nor is the density--zero nature of the prime denominators the
obstruction in itself: Bloom \cite{BloomShifted} represents every positive rational by distinct unit
fractions drawn from the (equally sparse) shifted primes $\{p-1\}$. What resists is the rank--one
rigidity --- the demand that the representation \emph{factor} as a product of two prime--reciprocal
sums --- which the barrier (Theorem~\ref{thm:barrier}) shows costs at least $59$ primes, hence of order
$870$ semiprime terms against Johnson's $47$.

Our starting point is an identification of \#307 with a question about the \emph{arithmetic
derivative} $n\mapsto n'$, the unique map with $0'=1'=0$, $p'=1$ for primes $p$, and the Leibniz
rule $(ab)'=a'b+ab'$. (The map goes back to J.~Mingot Shelly (1911) and appeared as a 1950 Putnam
problem; Barbeau \cite{Barbeau1961} gave the first systematic study, with the modern account in
Ufnarovski--\AA hlander \cite{UfnarovskiAhlander2003} and OEIS \cite{OEIS_A003415}.) For
$n=\prod_i p_i^{e_i}$ one has $n'=n\sum_i e_i/p_i$. We show (Section~\ref{sec:bridge}) that \#307
is \emph{equivalent} to the existence of a two--cycle of $n\mapsto n'$ that is not a fixed point.
The two--cycle reformulation itself is not new: it was stated ``in disguise'' by Seva
\cite{Seva323194} in 2019 (via the map $Q\mapsto Q\sum_{p\mid Q}1/p$, which agrees with the
arithmetic derivative on squarefree integers, where the question lives), and obtained independently
by Bado \cite{Bado2026}; see the note in Section~\ref{sec:prior}. Neither, however, identifies the
map with the \emph{arithmetic derivative} of Barbeau \cite{Barbeau1961} and Ufnarovski--\AA hlander
\cite{UfnarovskiAhlander2003}. That identification --- and the resulting two--way bridge to the
$n''=n$ literature (importing the Ufnarovski--\AA hlander conjecture, exporting our bounds) --- is the
starting contribution of this note; remarkably, Barbeau authored foundational papers on both sides
(1961 and 1976 \cite{Barbeau1976}).

The bridge is productive in both directions.
\begin{itemize}[leftmargin=1.4em]
\item \emph{Importing} from the derivative literature: the Ufnarovski--\AA hlander theorem that any
two--cycle has squarefree members with disjoint supports \cite{UfnarovskiAhlander2003} makes the
``finite sets of distinct primes'' hypothesis of \#307 automatic, and links \#307 to the
Ufnarovski--\AA hlander conjecture that all derivative cycles have period one.
\item \emph{Exporting} to that literature: our barrier ($\min(\prod P,\prod Q)\ge 2.09\times
10^{56}$, Theorem~\ref{thm:barrier}) improves Kovi\v{c}'s lower bound $\approx 3.2\times10^9$
\cite{Kovic2012} on the odd member of an $n''=n$ solution by roughly $47$ orders of magnitude.
\end{itemize}

\paragraph{Results and organisation.}
Section~\ref{sec:rigidity} proves the Rigidity Lemma (Theorem~\ref{thm:rigidity}): the fraction
$\sum_{p\in S}1/p$ is automatically in lowest terms, which forces, for any solution, the exact
identities $N_P=D_Q$, $N_Q=D_P$, disjointness $P\cap Q=\varnothing$, and $D_Q=s\,D_P$ where
$s=\sum_{p\in P}1/p$. The disjointness core was anticipated by R.~Israel \cite{MO320838}; the full
rigidity is what powers everything later. Section~\ref{sec:bridge} establishes the bridge.
Section~\ref{sec:barrier} proves the barrier, the parity dichotomy, and that longer cycles are
strictly harder (so two--cycles are extremal). Section~\ref{sec:anatomy} gives a local congruence
anatomy, a heuristic model, and a proof (via Erd\H{o}s--Wintner) that the relevant abundance density
is a genuine constant. Section~\ref{sec:ff} adds a function--field lens: over $\mathbb{F}_q[t]$ the
derivative has \emph{no} cycles, so the obstruction in $\mathbb{Z}$ is purely archimedean.
Section~\ref{sec:frame} develops the Frame Rule and divisor--trick,
and---crucially---records the exact identity that renders the natural constructive objective
circular. Section~\ref{sec:breeder} recasts the rule as a bilinear breeder, records a conditional
reduction $(\mathrm{H})\Rightarrow$ \textsc{yes} to an explicit prime--density hypothesis, and shows
quantitatively why $(\mathrm{H})$ does not amount to a feasible finite search.
Section~\ref{sec:computations} documents the computations, and Section~\ref{sec:lean} the
Lean~4 formalisation. We are explicit throughout about what is proved, what is heuristic, and what
is borrowed; the problem remains open.

\paragraph{Notation.}
For a finite set $S$ of distinct primes put
\[
D_S \coloneqq \prod_{p\in S} p, \qquad N_S \coloneqq \sum_{p\in S} \frac{D_S}{p},
\qquad\text{so}\qquad \sum_{p\in S}\frac1p = \frac{N_S}{D_S}.
\]
We call $N_S$ the \emph{cofactor sum}. For a solution we write $a=D_P$, $b=D_Q$,
$s=N_P/D_P=\sum_{p\in P}1/p$ and $t=N_Q/D_Q=\sum_{q\in Q}1/q$, so $st=1$, and (without loss of
generality) $s>1>t$. We write $\omega(n)$ for the number of distinct primes dividing $n$.

\section{The Rigidity Lemma}\label{sec:rigidity}

\begin{lemma}[Rigidity / lowest terms]\label{thm:rigidity}
For a finite set $S$ of distinct primes, $\gcd(N_S,D_S)=1$. Equivalently, $\sum_{p\in S}1/p$ is
already in lowest terms, and $v_p\!\bigl(\sum_{r\in S}1/r\bigr)=-1$ for every $p\in S$.
\end{lemma}

\begin{proof}
The prime divisors of $D_S=\prod_{r\in S}r$ are exactly the elements of $S$, so it suffices to show
$p\nmid N_S$ for each $p\in S$. Fix $p\in S$ and reduce $N_S=\sum_{r\in S}D_S/r$ modulo $p$. For
$r\ne p$ the term $D_S/r=\prod_{u\in S\setminus\{r\}}u$ still contains the factor $p$, so
$p\mid D_S/r$. The single surviving term is $D_S/p=\prod_{u\in S\setminus\{p\}}u$, whence
\[
N_S \equiv \prod_{u\in S\setminus\{p\}} u \pmod p .
\]
Each $u$ in this product is a prime distinct from $p$, so $p\nmid u$; as $p$ is prime it does not
divide the product. Thus $N_S\not\equiv0\pmod p$, i.e.\ $p\nmid N_S$. As $p$ ranged over all prime
divisors of $D_S$, $\gcd(N_S,D_S)=1$. The valuation statement is the same computation:
$v_p(N_S/D_S)=v_p(N_S)-v_p(D_S)=0-1=-1$.
\end{proof}

\begin{remark}
Distinctness and primality are both essential: $v_p$ of the cofactor of a repeated prime need not be
$-1$, and the final step uses Euclid's lemma. The case $|S|=1$ is consistent ($N_{\{p\}}=1$).
\end{remark}

\begin{theorem}[Solution structure]\label{thm:structure}
Suppose $P,Q$ are finite sets of distinct primes with $\dfrac{N_P}{D_P}\cdot\dfrac{N_Q}{D_Q}=1$.
Then
\[
N_P=D_Q,\qquad N_Q=D_P,\qquad P\cap Q=\varnothing,
\]
$N_P$ is squarefree with prime support exactly $Q$ (and $N_Q$ squarefree with support $P$), and
$D_Q=s\,D_P$ where $s=N_P/D_P$.
\end{theorem}

\begin{proof}
Cross-multiplying, $N_PN_Q=D_PD_Q$. By Lemma~\ref{thm:rigidity} both $N_P/D_P$ and $N_Q/D_Q$ are in
lowest terms. Since $N_Q/D_Q=D_P/N_P$ and the lowest--terms representative of a positive rational is
unique, comparing $N_Q/D_Q$ with $D_P/N_P$ gives $N_Q=D_P$ and $D_Q=N_P$.

\emph{Disjointness} (the Israel mechanism \cite{MO320838}). If $r\in P\cap Q$ then $r\mid D_P$ and,
by $N_P=D_Q=\prod Q$, also $r\mid N_P$; but $\gcd(N_P,D_P)=1$ forces $r\nmid N_P$, a contradiction.
(Equivalently $v_r(s)=v_r(t)=-1$, so $v_r(st)=-2\neq0=v_r(1)$.) Hence $P\cap Q=\varnothing$.

Finally $N_P=D_Q=\prod_{q\in Q}q$ is a product of distinct primes, hence squarefree with prime
support $Q$; symmetrically for $N_Q$. And $s\,D_P=(N_P/D_P)D_P=N_P=D_Q$.
\end{proof}

The disjointness in Theorem~\ref{thm:structure} means a solution is genuinely a partition of
$U\coloneqq P\cup Q$, and $D_PD_Q=\prod_{p\in U}p$. The valuation form $v_p(\sum 1/p_i)=-1$ underlying
disjointness is due to R.~Israel \cite{MO320838}, and appears also in Seva \cite{Seva323194}; the
cross--matching form ($N_P=D_Q$, $N_Q=D_P$) was obtained independently and concurrently by Bado
\cite[Thm.~A]{Bado2026} (Section~\ref{sec:prior}).

\section{The bridge to the arithmetic derivative}\label{sec:bridge}

\begin{lemma}\label{lem:der-cofactor}
For squarefree $m$, the arithmetic derivative equals the cofactor sum: $m'=\sum_{p\mid m}m/p=N_{\{p:\,p\mid m\}}$. Moreover $\gcd(m,m')=1$.
\end{lemma}

\begin{proof}
For $m=\prod_{p\mid m}p$ squarefree, $m'=m\sum_{p\mid m}1/p=\sum_{p\mid m}m/p$, which is exactly the
cofactor sum of the prime set of $m$. Coprimality is Lemma~\ref{thm:rigidity}.
\end{proof}

\begin{theorem}[Bridge]\label{thm:bridge}
The following are equivalent.
\begin{enumerate}[label=\textup{(\arabic*)},leftmargin=2.2em]
\item Erd\H{o}s \#307 has a solution: finite prime sets $P,Q$ with
$\bigl(\sum_{p\in P}1/p\bigr)\bigl(\sum_{q\in Q}1/q\bigr)=1$.
\item The arithmetic derivative has a two--cycle: integers $a\ne b$ with $a'=b$ and $b'=a$.
\item The equation $n''=n$ has a solution $n$ that is not a fixed point (i.e.\ $n'\ne n$, so $n$ is
not of the form $p^{p}$).
\end{enumerate}
\textup{(}The implication $(2)\Rightarrow(1)$ uses the Ufnarovski--\AA hlander squarefreeness
theorem for cycles; all other implications, and everything used in the barrier of
Section~\textup{\ref{sec:barrier}}, are unconditional --- see Remark~\textup{\ref{rem:cited}}.\textup{)}
\end{theorem}

\begin{proof}
$(1)\Rightarrow(2)$. Let $a=D_P,b=D_Q$. By Lemma~\ref{lem:der-cofactor},
$a'=N_P$ and $b'=N_Q$, and by Theorem~\ref{thm:structure}, $N_P=D_Q=b$ and $N_Q=D_P=a$. Thus
$a'=b$, $b'=a$, and $a\ne b$ since $s\ne1$ (a prime--reciprocal sum is never $1$, being
$N/D$ in lowest terms with $D>1$).

$(2)\Rightarrow(1)$. Let $a'=b,b'=a$ with $a\ne b$. By the Ufnarovski--\AA hlander analysis of
two--cycles (cf.\ \cite{UfnarovskiAhlander2003,Kovic2012}), both $a$ and $b$ are squarefree; their
prime supports $P,Q$ are then disjoint since $\gcd(a,b)=\gcd(a,a')=1$. By
Lemma~\ref{lem:der-cofactor}, $a'=N_P$ and $b'=N_Q$; the
cycle gives $N_P=b=D_Q$ and $N_Q=a=D_P$, so
$\bigl(\sum_{p\in P}1/p\bigr)\bigl(\sum_{q\in Q}1/q\bigr)=\frac{N_P}{D_P}\frac{N_Q}{D_Q}
=\frac{D_Q}{D_P}\frac{D_P}{D_Q}=1$.

$(2)\Leftrightarrow(3)$. A two--cycle $a\ne b$ gives $a''=a$ with $a'=b\ne a$, so $a$ is a non--fixed
solution of $n''=n$. Conversely a solution of $n''=n$ with $n'\ne n$ yields the two--cycle
$(n,n')$. The fixed points of $n\mapsto n'$ are exactly the numbers $p^{p}$
\cite{Barbeau1961,UfnarovskiAhlander2003}.
\end{proof}

\begin{remark}[Status of the cited input]\label{rem:cited}
Part $(2)\Rightarrow(1)$ uses that a derivative two--cycle has squarefree, support--disjoint
members. Squarefreeness is the content of the Ufnarovski--\AA hlander analysis of cycles (a
non--squarefree $n$ has $p\mid\gcd(n,n')$ for some $p$, and along $n\mapsto n'$ such common factors
propagate, precluding return); disjointness then follows from $\gcd(a,a')=1$. We use it as a cited
theorem. Even \emph{without} it, the implication $(1)\Rightarrow(2)$ and the entire barrier below are
unconditional, because Rigidity already delivers squarefree, disjoint $P,Q$ from a \#307 solution.
\end{remark}

\begin{corollary}\label{cor:ua}
If the Ufnarovski--\AA hlander conjecture holds in the weak form ``every cycle of $n\mapsto n'$ has
period one,'' then Erd\H{o}s \#307 has answer \textsc{no}. Conversely, a \textsc{yes} to \#307
exhibits a period--two cycle and refutes that statement.
\end{corollary}

\section{The barrier}\label{sec:barrier}

\begin{lemma}[Strict AM--GM]\label{lem:amgm}
If $s,t>0$ with $st=1$ and $s\ne1$, then $s+t>2$.
\end{lemma}
\begin{proof}
$s+t-2=s+1/s-2=(s-1)^2/s>0$ since $s\ne1$.
\end{proof}

\begin{lemma}[The number $59$, after Rosen \cite{MO320838}]\label{lem:59}
Any finite set $U$ of distinct primes with $\sum_{p\in U}1/p>2$ has $|U|\ge 59$. Moreover, among
$k$--element prime sets, $\sum 1/p$ is maximised by the $k$ smallest primes.
\end{lemma}
\begin{proof}
Replacing any prime of $U$ by a smaller prime not in $U$ increases $\sum1/p$, so the maximum over
$k$--element sets is at the first $k$ primes. The reciprocal sum of the first $58$ primes is
$1.99874\ldots<2$ and of the first $59$ primes is $2.00235\ldots>2$, so a sum exceeding $2$ needs at
least $59$ primes.
\end{proof}

\begin{theorem}[Barrier]\label{thm:barrier}
For any solution $(P,Q)$ of \#307, with $U=P\cup Q$:
\begin{enumerate}[label=\textup{(\alph*)},leftmargin=2.2em]
\item $|U|\ge 59$ and $\displaystyle\prod_{p\in U}p \ge \Pi_{59}\coloneqq\prod_{i=1}^{59}p_i
\approx 8.77\times10^{112}$;
\item $\displaystyle\min(\textstyle\prod P,\prod Q)\ \ge\ \sqrt{\Pi_{59}/T_{59}} \ =\
2.0929\ldots\times10^{56}$, where $T_{59}=\sum_{i=1}^{59}1/p_i=2.00235\ldots$. In particular both
$\prod P$ and $\prod Q$ exceed $2.09\times10^{56}$.
\end{enumerate}
Consequently every $n''=n$ solution has both members $>2.09\times10^{56}$, improving Kovi\v{c}'s
bound $\approx 3.2\times10^9$ \cite{Kovic2012} by about $47$ orders of magnitude.
\end{theorem}

\begin{proof}
(a) By $st=1$, $s\ne1$ and Lemma~\ref{lem:amgm}, $\sum_{p\in U}1/p=s+t>2$, so $|U|\ge 59$ by
Lemma~\ref{lem:59}; since $P,Q$ are disjoint (Theorem~\ref{thm:structure}),
$\prod_{p\in U}p=D_PD_Q$, and a product of $\ge 59$ distinct primes is at least the product of the
$59$ smallest, $\Pi_{59}$.

(b) Take $D_P\le D_Q$ to be the minimal side. By Theorem~\ref{thm:structure}, $D_Q=s\,D_P$ with
$s>1$, hence
\[
s\,D_P^2 \;=\; D_P\,D_Q \;=\; \prod_{p\in U}p \;=\!:\, R,
\qquad\text{so}\qquad D_P^2=\frac{R}{s}.
\]
Since $s\le s+t=T(U)\coloneqq\sum_{p\in U}1/p$ (as $t>0$),
\[
D_P^2=\frac{R}{s}\ \ge\ \frac{R}{T(U)} .
\]
It remains to bound $R/T(U)$ below over all admissible $U$. Among $k$--element prime sets the
quotient $\bigl(\prod U\bigr)\big/\bigl(\sum_{p\in U}1/p\bigr)$ is minimised by the $k$ smallest
primes (smallest numerator and, by Lemma~\ref{lem:59}, largest denominator simultaneously), and
passing from $59$ to more primes only increases it (the product grows multiplicatively while the
reciprocal sum grows by a single small term). Hence the global minimum over admissible $U$ is at
$U=\{p_1,\dots,p_{59}\}$, giving $R/T(U)\ge \Pi_{59}/T_{59}$. Therefore
\[
D_P^2 \ \ge\ \frac{\Pi_{59}}{T_{59}} \;=\; \frac{\Pi_{59}^2}{N_{59}}\;=\;4.3806\ldots\times10^{112},
\qquad D_P\ \ge\ 2.0929\ldots\times10^{56},
\]
where $N_{59}=T_{59}\,\Pi_{59}$ is the integer cofactor sum of the first $59$ primes.
\end{proof}

\begin{remark}[Why the \emph{minimum} is bounded, not only the maximum]
The product identity $D_PD_Q\ge\Pi_{59}$ alone gives only $\max(\prod P,\prod Q)\ge
\sqrt{\Pi_{59}}\approx2.96\times10^{56}$. The stronger statement~(b), a lower bound on the
\emph{smaller} side, genuinely uses rigidity: it is the identity $D_Q=s\,D_P$ (equivalently
$D_P^2=R/s$) together with $s\le T(U)$ that converts the product bound into a bound on $D_P$. We
verified numerically that $\min_U R(U)/T(U)$ is attained exactly at the first $59$ primes (the
minimising configuration is unique), so the constant $2.09\times10^{56}$ is sharp for this method.
\end{remark}

\begin{proposition}[Closing the $59$--prime level]\label{prop:close59}
No solution of \#307 has $|P\cup Q|=59$. Hence $|P\cup Q|\ge60$,
$\prod_{p\in U}p\ \ge\ \Pi_{60}>2.46\times10^{115}$, and
\[ \min\bigl(\textstyle\prod P,\prod Q\bigr)\ \ge\ \sqrt{\Pi_{60}/T_{60}}\ =\
3.5053\ldots\times10^{57}. \]
\end{proposition}

\begin{proof}[Proof (finite verification)]
Let $U=P\cup Q$ with $|U|=59$; as in Theorem~\ref{thm:barrier}(a), $T(U)>2$. The candidate list is
finite and complete: \textup{(i)} every prime $p\le167$ lies in $U$ --- omitting $p$ caps $T(U)$ at
$T_{60}-1/p<2$, since the $59$ largest prime reciprocals avoiding $p$ are those of the first $60$
primes less $p$, and $1/p>T_{60}-2=0.00590\ldots$ exactly for $p\le167$; \textup{(ii)} every element
$p'\ge281$ satisfies $T_{58}+1/p'>2$ (the other $58$ reciprocals sum to at most $T_{58}$), so
$p'\le787$. Thus $U$ consists of the $39$ primes up to $167$ plus twenty primes from
$(167,787]$, and exactly $49{,}961$ such sets have $T(U)>2$. For each, a solution supported on $U$
forces both $N'+2N=(a+b)^2$ and $N'-2N=(a-b)^2$ to be perfect squares, where $N=\prod U$ and
$N'=a^2+b^2$ (rigidity; the $k=0$ case of Proposition~\ref{prop:pairform}) --- one integer test that
simultaneously excludes all $2^{58}$ splits of $U$. Exact--integer verification shows $N'+2N$ is a
non--square for \emph{every} one of the $49{,}961$ candidates, by two independent implementations
with different enumeration orders and prunings (\texttt{code/close59.py} and
\texttt{hunt/close59.rs} in the repository). Hence no $59$--element support closes, and
$|U|\ge60$. The bound on the minimal side repeats the proof of Theorem~\ref{thm:barrier}(b) with the
minimum of $R/T$ now over $|U|\ge60$, attained at the first $60$ primes:
$D_P^2\ge\Pi_{60}/T_{60}=\Pi_{60}^2/N_{60}$, where $N_{60}$ is the cofactor sum of the first $60$
primes.
\end{proof}

\begin{corollary}[The coprime variant inherits the closed level]\label{cor:coprime60}
Consider the weakened variant in which the elements of $P$ and $Q$ are merely pairwise coprime
integers $\ge2$ \textup{(}the open form of the variant, $1\notin P\cup Q$\textup{)}. Any solution
has $|P\cup Q|\ge 60$ elements, and $\min(\prod P,\prod Q)>3.50\times10^{57}$.
\end{corollary}
\begin{proof}
Rigidity holds verbatim for pairwise--coprime families (the cofactor--sum argument uses only
coprimality), so again $T(U)=s+t>2$. The least prime factors $\ell(e)$ of the elements are distinct,
and $\sum_e 1/e\le\sum_e1/\ell(e)$. Suppose $|U|=59$ and some element $e$ is not prime, so
$e\ge\ell(e)^2$. If $\ell(e)\le277$, the loss is
$1/\ell-1/e\ge(\ell-1)/\ell^2\ge276/277^2=0.00359\ldots$, exceeding the maximal available slack
$T_{59}-2=0.00235\ldots$; if $\ell(e)\ge281$, then $e\ge281^2$ while the other $58$ reciprocals sum
to at most $T_{58}$, giving $T(U)\le T_{58}+281^{-2}=1.99875\ldots<2$. Either way $T(U)<2$, a
contradiction: at $59$ elements every element is prime, and Proposition~\ref{prop:close59} applies
verbatim. The product bound transfers because $\prod_e e\ge\prod_e\ell(e)$ and
$T(U)\le\sum_e1/\ell(e)$, so the extremal ratio bound at the first $60$ primes still applies.
\end{proof}

\begin{remark}[Verification tier, and why the ladder stops here]\label{rem:tier60}
Theorem~\ref{thm:barrier} is machine--checked in Lean~4; Proposition~\ref{prop:close59} is a
computer--assisted finite verification (two independent exact--integer programs), not yet
formalised. The level cannot be climbed further by finite means: already at $|U|=60$ the family
$U=\{p_1,\dots,p_{59}\}\cup\{q\}$ qualifies for \emph{every} prime $q>277$, and on it the two square
conditions read $x^2=Aq+\Pi_{59}$, $y^2=Bq+\Pi_{59}$ with $A,B=N_{59}\pm2\Pi_{59}$, i.e.\ the
Pell--type system $Bx^2-Ay^2=-4\Pi_{59}^2$ with $113$--digit coefficients. We verified that no
accessible modulus obstructs it (mod $32$, and all odd $\ell<4000$: the conditions pin
$q\equiv5\pmod 8$ and thin each $\ell\mid\Pi_{59}$ to about half of its residue classes and each
$\ell\nmid\Pi_{59}$ to about a quarter, but never to zero) --- consistent with the local--solubility
theme of Section~\ref{sec:anatomy}: the residue of the problem is again global, not local. A
genuine obstruction would have to come from the primes dividing $A$ or $B$, which is
factoring--hard at this size. The tail family also exposes the self--similarity of the ladder:
writing $a=qm$, so that $m\,b=\Pi_{59}$ partitions the first $59$ primes, its two closing
conditions collapse to the exact equation
\[ D(m)\,D(b)+m^2\;=\;\Pi_{59} \]
together with the divisibility $m\mid D(b)$ (then $q=D(b)/m$, required prime) --- verified
exhaustively on a $12$--prime toy universe. Equivalently, the deficit $1-\sigma(m)\sigma(b)$ must
equal $m^2/\Pi_{59}$ \emph{exactly}: the circularity $\Delta_0=M_0N(1-s_{M_0}s_N)$ of
Section~\ref{sec:frame} in its purest form, with the deviation forced to be the perfect square
$m^2$. Closing even this one family is thus a $2^{58}$--fold exact--coincidence search of the same
type as \#307 itself (expected count ${\sim}2^{58}/\Pi_{59}\approx10^{-95}$): each rung of the
ladder past $60$ contains the whole problem.
\end{remark}

\begin{theorem}[Parity dichotomy]\label{thm:parity}
Let $a'=b,b'=a$ be a derivative two--cycle (so $a,b$ squarefree, $\gcd(a,b)=1$). Then exactly one of
the following holds.
\begin{enumerate}[label=\textup{(\roman*)},leftmargin=2.2em]
\item \emph{Mixed parity.} Exactly one of $a,b$ is even; say $a=2k$ with $k$ odd. Then $b=a'$ is odd
and $\omega(b)$ is even.
\item \emph{All odd.} Both $a,b$ are odd. Then $P\cup Q$ avoids $2$, so reaching reciprocal sum $>2$
needs $\ge 1412$ odd primes, and $\min(a,b)>10^{2519}$.
\end{enumerate}
\end{theorem}

\begin{proof}
They cannot both be even, as $\gcd(a,b)=\gcd(a,a')=1$. In case~(i), write $a=2k$, $k$ odd squarefree;
$b=a'=k+2k'$ is odd. For odd squarefree $b=\prod_{i}q_i$ one has
$b'=\sum_i b/q_i\equiv \omega(b)\pmod 2$ (each $b/q_i$ is odd). Since $b'=a$ is even,
$\omega(b)\equiv0\pmod2$. In case~(ii) the union of supports omits $2$; the reciprocal sum of the
first $1411$ odd primes is $<2$ while that of the first $1412$ exceeds $2$, so $|P\cup Q|\ge 1412$
and $\prod_{p\in P\cup Q}p\ge 10^{5040}$ (product of the $1412$ smallest odd primes, by direct
computation). Applying the same rigidity/ratio argument as in Theorem~\ref{thm:barrier}(b) --- now
with $2\notin U$ --- gives $\min(a,b)\ge\sqrt{R/T(U)}>10^{2519}$. The all--odd case is therefore
vastly more constrained than the mixed case and, heuristically, empty.
\end{proof}

\begin{proposition}[Longer cycles are harder; two--cycles are extremal]\label{prop:kcycles}
Every member of a $k$--cycle of the arithmetic derivative is squarefree, and for $k\le3$ the members
are pairwise coprime (each pair is consecutive). Their masses $s(a_i)=a_{i+1}/a_i$ multiply to $1$,
so by AM--GM $\sum_i s(a_i)\ge k$; for $k\le3$ this is the union's reciprocal sum, whence the support
contains at least $m_k$ primes, where $m_k$ is least with $\sum_{i\le m_k}1/p_i\ge k$. Now $m_2=59$
and $m_3=361{,}139$ (last prime $5{,}195{,}977$), so a three--cycle has at least $361{,}139$ primes in
its support and product exceeding $10^{2.26\times10^{6}}$ (hence its largest member exceeds
$10^{752{,}194}$); by Mertens $m_k$ grows doubly exponentially. Two--cycles are thus the extremal ---
and only practically relevant --- case. (For $k\ge4$ non--consecutive members may share primes; the
bound then holds for the multiset reciprocal sum.)
\end{proposition}

\section{Local anatomy and a heuristic model}\label{sec:anatomy}

\begin{lemma}[Local anatomy]\label{lem:anatomy}
Let $a$ be squarefree and $q$ a prime with $q\nmid a$. Then $q\mid a'$ if and only if
$\sum_{p\mid a}p^{-1}\equiv 0\pmod q$, where $p^{-1}$ is the inverse of $p$ modulo $q$. If $q\mid a$,
then $q\nmid a'$.
\end{lemma}
\begin{proof}
$a'=a\sum_{p\mid a}1/p$. If $q\nmid a$ then $q\mid a'$ iff $q\mid a\sum_{p\mid a}p^{-1}$ iff
$\sum_{p\mid a}p^{-1}\equiv0\pmod q$. If $q\mid a$, then by Lemma~\ref{thm:rigidity}
applied to the prime set of $a$, $q\nmid a'$.
\end{proof}

Thus the cycle equation $a'=b$, $b'=a$ factors, prime by prime, into
$\omega(a)+\omega(b)$ local congruences of exactly the shape that governs amicable--pair heuristics:
for each $q\mid b$ one needs $\sum_{p\mid a}p^{-1}\equiv0\pmod q$, and symmetrically. These local
conditions, however, are individually \emph{empty of obstruction}:

\begin{proposition}[No single--modulus congruence obstruction]\label{prop:nolocal}
For every modulus $m$ with $2\le m\le210$, the reduced cross--match system $N_P\equiv D_Q$,
$N_Q\equiv D_P\pmod m$ is realisable by disjoint squarefree prime sets $P,Q$, in both structural
branches: $\gcd(m,D_PD_Q)=1$, and the branch in which a prime divisor of $m$ lies in $P\cup Q$. Hence
no individual modulus in this range can rule out a solution of \#307.
\end{proposition}
\begin{proof}[Construction]
Write $\Sigma_S\equiv\sum_{p\in S}p^{-1}\pmod m$, so $N_S\equiv D_S\Sigma_S$. The reduced system is
$\Sigma_P\equiv D_Q D_P^{-1}$, $\Sigma_Q\equiv D_P D_Q^{-1}\pmod m$; its only coupling is
$\Sigma_P\Sigma_Q\equiv1$, automatic, with $D_P,D_Q$ otherwise free. By Dirichlet each unit class mod
$m$ contains primes, fixing $D_P,D_Q$; appending a prime $r\equiv1\pmod m$ leaves $D_S$ fixed while
sending $N_S\mapsto rN_S+D_S\equiv N_S+D_S$, i.e.\ incrementing $\Sigma_S$ by $1$, so (when
$\gcd(D_S,m)=1$) every target class is reached with disjoint sets. In the branch $\ell\mid m$,
$\ell\in P$, reduction mod $\ell$ collapses the system to $N_Q\equiv D_P\equiv0\pmod\ell$, met by
inverse--paired residues in $Q$; composite $m$ needs no separate gluing, the increment acting directly
mod $m$. We verified both branches for all $2\le m\le210$ (for prime and prime--power moduli the
inside branch requires the directed construction above --- seeding a prime divisor of $m$ into $P$ ---
a blind search missing these as a search--power artefact, not an arithmetic obstruction).
\end{proof}
\begin{remark}[Precise scope]
This excludes only a \emph{single fixed--modulus} congruence disproof. The witnesses are
$m$--dependent (the modulus--$6$ witness already fails modulo $210$), and a single $(P,Q)$ solving the
system for \emph{all} $m$ simultaneously is, for bounded integers, equivalent to the integer identities
$N_P=D_Q$, $N_Q=D_P$ --- that is, to \#307 itself --- so per--modulus solvability is \emph{no} evidence
for existence. Nor does it exclude other finite or local negative arguments (size/Mertens bounds,
density, magnitude--congruence hybrids, or a Hasse--type failure across infinitely many moduli);
indeed our own $|U|\ge59$ and $2.09\times10^{56}$ barriers are finite--data results fully consistent
with it. Together with Theorem~\ref{thm:ff} --- every non--archimedean shadow trivially negative for
degree reasons --- it brackets the problem: no single--modulus obstruction over $\mathbb{Z}$, fatal
non--archimedean analogues over $\mathbb{F}_q[t]$, the difficulty living at the archimedean place.
\end{remark}

\paragraph{A first--order count.}
Model a candidate squarefree $a$ as ``$P$--abundant'' ($\sum_{p\mid a}1/p>1$) and ask for
$b=a'$ to itself be $Q$--abundant with $\sum_{q\mid b}1/q=1/s$. Two empirical inputs (Section
\ref{sec:computations}) calibrate this:
\begin{itemize}[leftmargin=1.4em]
\item the density of $P$--abundant $n$ is $\delta=0.0420$ (and $\approx0.0176$ among squarefree
$n$) --- a theorem, not merely an observation (Proposition~\ref{prop:density});
\item for a qualifying squarefree $a$, $b=a'$ is squarefree about $95\%$ of the time, but its mass
$\sum_{q\mid b}1/q$ has mean $\approx0.047$ against the required $\approx0.934$ --- the empirical tail
probability $\Pr[\sum_{q\mid b}1/q\ge 1/s]$ is $0$ in $40{,}920$ samples at these scales
(Section~\ref{sec:computations}).
\end{itemize}
The gap is structural: at small scale a $P$--abundant $a$ hogs the cheap mass--carrying small primes,
and $b=a'$, being coprime to $a$, is both too small and too coprime to assemble reciprocal mass
$1/s$. The extremal case is stark: for the primorial $a=\prod_{i\le k}p_i$, though $\sum_{p\mid a}1/p$
rises through $2$ at $k=59$, the cofactor sum $a'$ is essentially prime, with $\sum_{q\mid a'}1/q<0.2$
for every $k<100$ (and only $0.0014$ at $k=59$): the family that \emph{maximises} the forward mass is
the most starved on the return. The first scale at which $b$ could plausibly carry the required mass is of the order of the
barrier ($\sim10^{56}$); this dynamical mass--gap is \emph{consistent with} the proved barrier rather
than an independent derivation of it (the empirical tail probability $0$ at small scale is equally
consistent with \textsc{no} and with \textsc{yes}--beyond--the--barrier). A naive ``expected number
of solutions'' built from the local congruences (one factor $\sim1/q$ per condition) diverges like a
constant times $\log X$; this is a divergent naive count of exactly the type known to be unreliable,
so we treat it only as very weak, non--rigorous suggestive evidence that solutions could exist at or
beyond the barrier.

\begin{remark}[Calibration]
We present the count of the previous paragraph as a heuristic, not a theorem. Heuristics of
amicable--pair type are known to be delicate; we report it as weak evidence, consistent with the
divergence of the refined expected count, that the answer to \#307 is plausibly \textsc{yes}
(and hence that the period--one form of the Ufnarovski--\AA hlander conjecture is plausibly false).
\end{remark}

\begin{remark}[A cautionary refinement]\label{rem:refined}
One can try to sharpen this into a quantitative rate using a Kubilius/Erd\H{o}s--Wintner model,
stratified over which small primes lie in the support of $a$ versus remain available to $b$. Two
lessons emerge, both methodological. First, a plain Monte--Carlo estimate \emph{undercounts} the rate
by some $10^{61}$: the expectation is dominated by rare, balanced splits of the small primes that
sampling almost never sees, so a naive estimate of this kind is misleadingly small and only an
explicit stratified sum corrects it. Second --- and this is why we record no numerical prediction ---
the stratified estimate is itself unstable: refining the stratification from the first $8$ to the
first $14$ primes drops the rate by some $17$ orders of magnitude and pushes the implied first
occurrence from $\sim10^{10^{38}}$ out past $\sim10^{10^{55}}$, with no sign of convergence. The only
robust conclusion is negative: every version places a hypothetical first solution \emph{doubly
exponentially} beyond any conceivable search, consistent with the constructive deficit of
Section~\ref{sec:breeder}; the model does not reliably decide existence in either direction.
\end{remark}

\begin{proposition}[The abundance density is a theorem]\label{prop:density}
The additive function $f(n)=\sum_{p\mid n}1/p$ satisfies the Erd\H{o}s--Wintner three--series
criterion \cite{Tenenbaum} (here $\sum 1/p^2$ and $\sum 1/p^3$ converge), so $f$ has a limiting
distribution and the
density $\delta=\lim_{x\to\infty}\tfrac1x\#\{n\le x:f(n)>1\}$ exists, equal to $\Pr\bigl[\sum_p
X_p/p>1\bigr]$ for independent $X_p\sim\mathrm{Bernoulli}(1/p)$. Exact convolution over primes
$\le10^5$ (with a Markov tail bound) gives $\delta=0.0420\pm0.001$, agreeing to four places with the
sieve to $2\times10^7$. As for the density of abundant integers, no closed form is expected.
\end{proposition}

\begin{remark}[Two dynamical probes: no second obstruction]\label{rem:dynamical}
Two measurements sharpen the picture beyond the first--order mass count, and both cut towards rather
than against a solution. \emph{(i) Image density.} A two--cycle needs \emph{both} members in the image
$D(\mathbb{N})$ (each is the other's derivative), so one might fear a second obstruction --- that the
high--mass numbers a cycle requires are seldom derivatives. The opposite holds: of the integers below
$10^7$, $68.3\%$ lie in $D(\mathbb{N})$, and the fraction \emph{rises} with mass, from $61.7\%$ at
$\sum_{p\mid n}1/p<0.2$ to $78.2\%$ at $\sum_{p\mid n}1/p>1$ and $80\%$ beyond. The high--mass numbers
are thus over--represented among derivatives; the second cycle condition is comfortably met.
\emph{(ii) The barrier is where the balanced mass--product reaches $1$.} Set
$r(n)=\bigl(\sum_{p\mid n}1/p\bigr)\bigl(\sum_{q\mid n'}1/q\bigr)=D^2(n)/n$, so $r=1$ is exactly the
cycle condition. Exhaustively over $n\le X$ the maximum of $r$ climbs $0.50\to0.62$ as $X$ runs
$10^3\to10^9$ (with no exact cycle, as the barrier demands), with no ceiling; the climb follows the
balanced split $r\approx(S(k)/2)^2$, $S(k)=\sum_{i\le k}1/p_i$, which reaches $1$ precisely when
$S(k)=2$, i.e.\ $k\approx59$. The dynamical trajectory thus reproduces the $|U|\ge59$ barrier from the
opposite direction. Neither probe decides existence; both \emph{remove} hypothetical obstructions,
leaving the exact archimedean closure as the sole gate.
\end{remark}

\begin{remark}[Why the analytic method is starved: $N(r)\le1$]\label{rem:nostructure}
The rank--one rigidity (Section~\ref{sec:status}) has a sharp consequence for the circle--method
machinery that resolves the \emph{unrestricted} Egyptian--fraction problem. For a rational $r$ let
$N(r)=\#\{\,P:\sum_{p\in P}1/p=r\,\}$ count the finite prime sets with reciprocal sum $r$. Then
\[N(r)\le1:\]
by the Rigidity Lemma (Theorem~\ref{thm:rigidity}) the sum $\sum_{p\in P}1/p=N_P/D_P$ is already in
lowest terms, so the reduced denominator of $r$ equals $D_P=\prod P$ and $P$ is exactly its set of
prime divisors --- uniquely determined ($N(r)=1$ precisely when that denominator is squarefree and the
numerator of $r$ is its cofactor sum). The circle method needs the opposite: a target hit by a positive
proportion of subsets, so a main term dominates. This is exactly what Bloom's shift to $\{p-1\}$
manufactures --- $p-1$ being composite and divisor--rich, a rational is the reciprocal sum of
\emph{many} subsets of $\{1/(p-1)\}$ (already among the first fourteen shifted primes the multiplicity
reaches $3$ and grows), whence \cite{BloomShifted} represents every positive rational there. Over the
primes themselves $N(r)\le1$ leaves nothing to average: the analytic engine that reaches $\{p-h\}$ is
provably starved on $\{p\}$, and this is the method--level reason \#307 --- the rank--one, exact--prime
case --- lies beyond it.
\end{remark}

\section{A function--field lens: the obstruction is archimedean}\label{sec:ff}

Define an arithmetic derivative on $\mathbb{F}_q[t]$ by $P'=1$ for monic irreducibles and the Leibniz
rule, so $f'=\sum_i e_i\,f/P_i$ for $f=\prod_i P_i^{e_i}$ (the exponents $e_i$ read modulo the
characteristic).

\begin{theorem}\label{thm:ff}
Over $\mathbb{F}_q[t]$ this derivative has no two--cycles and no nonzero fixed points. Indeed
$\deg f'\le\deg f-1$ for every nonconstant $f$. The analogues of the integer fixed points $p^{p}$ also
vanish: $(P^{p})'=p\,P^{p-1}P'=0$ in characteristic $p$. The same holds verbatim for every twisted
variant ($P'=u_P\in\mathbb{F}_q^\times$, Leibniz).
\end{theorem}
\begin{proof}
Each term $e_i\,f/P_i$ has degree $\deg f-\deg P_i\le\deg f-1$, so $\deg f'\le\deg f-1$. A two--cycle
$f'=g$, $g'=f$ would force $\deg g\le\deg f-1$ and $\deg f\le\deg g-1$, hence $\deg f\le\deg f-2$;
a fixed point $f'=f$ is excluded likewise. (Verified exhaustively over $\mathbb{F}_2$ to degree $7$,
$\mathbb{F}_3$ to degree $5$, and $\mathbb{F}_5$ to degree $4$; twisted variants over $\mathbb{F}_2$ and
$\mathbb{F}_5$ on small supports through $2.3\times10^{6}$ twist configurations, none closing.)
\end{proof}

\begin{corollary}[The function--field Erd\H{o}s--Barbeau equation is insoluble]\label{cor:ff307}
For every finite field $\mathbb{F}_q$, all nonempty finite sets $P,Q$ of monic irreducibles, and all
unit weights $u_\pi,v_\rho\in\mathbb{F}_q^\times$,
\[
\Bigl(\sum_{\pi\in P}\frac{u_\pi}{\pi}\Bigr)\Bigl(\sum_{\rho\in Q}\frac{v_\rho}{\rho}\Bigr)\ne1 .
\]
\end{corollary}
\begin{proof}
Write $f=\prod P$, $g=\prod Q$, $D_u(f)=\sum_\pi u_\pi f/\pi$; the equation reads $D_u(f)\,D_v(g)=fg$.
If some $\pi\in P\cap Q$, then $\pi^2\mid fg$ while $D_u(f)\equiv u_\pi f/\pi\not\equiv0$ and
$D_v(g)\equiv v_\pi g/\pi\not\equiv0\pmod\pi$, so $\pi\nmid D_u(f)D_v(g)$ --- impossible; hence
$\gcd(f,g)=1$. The same congruence gives $\gcd(f,D_u(f))=\gcd(g,D_v(g))=1$, so $g\mid D_u(f)$ and
$f\mid D_v(g)$, forcing $D_u(f)=c\,g$, $D_v(g)=c^{-1}f$ with $c\in\mathbb{F}_q^\times$: a twisted
two--cycle, excluded by the degree drop. The weighted product equation is thus insoluble over every
$\mathbb{F}_q(t)$; it \emph{is} soluble over the unit--rich number rings of
Theorem~\ref{thm:existence}, since any twisted cycle gives $(D_u(a)/a)(D_v(b)/b)=(b/a)(a/b)=1$; and over
$\mathbb{Q}$, where the admissible weights are $\pm1$, it is \#307 itself with the signed barrier of
Remark~\ref{rem:ordered}, open beyond $2.09\times10^{56}$. The three regimes --- non--archimedean
monotonicity (closed, negative), unit--broken monotonicity (closed, positive), ordered--archimedean
(open, barred) --- now have theorems on two of three doors.
\end{proof}

Over $\mathbb{Z}$, by contrast, $n'=n\sum_i e_i/p_i$ can exceed $n$ precisely because the archimedean
absolute value lets a sum of strictly smaller terms outgrow the original --- the ``mass $>1$''
phenomenon $\sum 1/p>1$, which has no degree--theoretic analogue. Theorem~\ref{thm:ff} thus localises
the entire existence question to the archimedean place: every non--archimedean shadow of \#307 is
trivially negative. This is consistent with (and partly explains) both the census finding of
Section~\ref{sec:computations} and Proposition~\ref{prop:nolocal} (no single modulus obstructs the
cross--match system in the tested range) --- the difficulty is not local. Any eventual proof must therefore distinguish $\mathbb{Z}$ from $\mathbb{F}_q[t]$; over the
latter the decisive tools are the degree and the Mason--Stothers/Wronskian inequality
\cite{Stothers}, neither of which has a $\mathbb{Z}$--analogue. The conjectural analogue --- $abc$ ---
does not fill the gap, and quantifiably so. Every cycle inherits, from $N'=a^2+b^2$ with $N=ab$, the
relation $(a-b)^2+4ab=(a+b)^2$ on coprime squarefree members; but this triple is $abc$--unexceptional,
of quality $q=\log(a+b)^2/\log\operatorname{rad}\!\bigl((a-b)\cdot ab\cdot(a+b)\bigr)\approx\tfrac12$
(median $0.543$, never above $0.74$ across $6\times10^4$ coprime squarefree pairs; the only $q>1$ cases
are degenerate smooth--gap pairs such as $(5,3)$, vanishingly far below the $10^{56}$ scale a solution
must occupy). Granting $abc$ in its strong explicit form would therefore yield no bound here: the
function--field collapse rests on an inequality genuinely \emph{absent} over $\mathbb{Z}$, not merely
unproven, and the size mismatch is by a factor of two in the exponent. \#307 lives in the
$abc$--comfortable zone.

\begin{remark}[The free grading, and what evaluation destroys]\label{rem:grading}
Theorem~\ref{thm:ff} has a structural reading. In the free Leibniz world --- the polynomial ring
$\mathbb{Z}[x_p]$ with $x_p'=1$ --- the derivative of a squarefree monomial of degree $n$ is the
elementary symmetric form $e_{n-1}$, homogeneous of degree $n-1$: the formal derivative strictly lowers
the grading, cycles are impossible, and over $\mathbb{F}_q[t]$ the degree realises this grading
faithfully. Over $\mathbb{Z}$ the evaluation $x_p\mapsto p$ destroys homogeneity --- and the data show it
destroys it completely: for random squarefree $a$ of fixed size, $\omega(a')$ is governed by the size of
$a'$ alone (Erd\H{o}s--Kac), flat in $\omega(a)$ within each parity class (measured: mean
$\omega(a')=2.7$--$3.0$ as $\omega(a)$ runs over $3,5,7$, and $3.6$--$4.0$ over $4,6,8$, the offset being
the parity anatomy of Theorem~\ref{thm:parity}). A two--cycle is therefore a doubly anti--typical event:
an additively assembled integer must carry multiplicative rank far in the Erd\H{o}s--Kac tail \emph{and}
on an exactly prescribed support --- the two independent prices whose product is the heuristic rarity of
Section~\ref{sec:anatomy}. The deeper additive--multiplicative content collapses, on inspection, into
the $S$--unit form of Proposition~\ref{prop:sunit}: dividing the cycle equation by $D_Q$ exhibits it as
an $n$--term unit equation --- one more coordinate system in which the same exactness reappears
unchanged.
\end{remark}

\begin{remark}[The obstruction is \emph{ordered}--archimedean]\label{rem:ordered}
The barrier is finer than ``archimedean'': its engine is AM--GM on real, \emph{positive} masses, i.e.\
the ordering of $\mathbb{R}$. Over $\mathbb{Z}[i]$ (a PID, so the Rigidity Lemma and cross--matching
survive verbatim, with cycles defined up to units), the masses $\sum1/\pi$ are \emph{complex}: phases can
align so that $\bigl|\sum1/\pi\bigr|>1$ with as few as three Gaussian primes (e.g.\
$\tfrac1{1+i}+\tfrac1{2+i}+\tfrac1{2-i}=1.3-0.5i$), and the $59$--prime bound has no analogue. The
ordered--archimedean thesis thus predicts that small derivative two--cycles \emph{may} exist over
$\mathbb{Z}[i]$ where they are barred over $\mathbb{Z}$ --- and the prediction is confirmed. Two facts
make the contrast precise.
\emph{First}, phases do not help over $\mathbb{Z}$: for any \emph{signed} derivative ($p'=\pm1$, Leibniz)
a two--cycle still satisfies $|s||t|=1$ with $|s|\le\sigma(a)$, $|t|\le\sigma(b)$ and disjoint supports
(the rigidity proof is unchanged), so $\sigma(a)+\sigma(b)\ge2$ and the full $59$--prime,
$2.09\times10^{56}$ barrier applies to every choice of signs: real phases cannot cancel.
\emph{Second}, fourth roots of unity suffice instantly: Leibniz derivatives on $\mathbb{Z}[i]$ with unit
prime--values ($\pi'\in\{\pm1,\pm i\}$) possess genuine two--cycles at norm ${\sim}10^4$. The smallest:
\[ a=(1+i)(2+7i)(6+11i)=-129-i,\qquad b=3\,(2+i)(1+2i)(6+5i)=-75+90i, \]
with $a'=b$ under the \emph{standard} values $\pi'=1$ on the primes of $a$, and $b'=a$ exactly under
$\pi'=i$ on $3$, $2+i$, $1+2i$ and $\pi'=1$ on $6+5i$ --- seven prime classes
($N(a)=16{,}642$, $N(b)=13{,}725$) against the fifty--nine that $\mathbb{Z}$ demands. Exhaustive search
finds exactly sixteen such cycles up to conjugation and direction having a member of norm
$\le2\times10^7$ with $\omega\le5$ on the enumerated side --- $\omega$--profiles $(3,4)$, $(4,4)$,
$(5,4)$, and even $(4,6)$ --- and the population grows steadily ($4$ below $3\times10^5$, $16$ below
$2\times10^7$), consistent with an infinite, slowly growing family; in six of the sixteen one side
requires no twist at all. All sixteen were verified exactly, by independent recomputation; no twisted
\emph{fixed} point ($a'=ua$) exists in the same range. The strict two--sided normalisation $\pi'=1$
(first--quadrant representatives) has no cycle up to units --- nor any fixed point, zero--derivative, or
unit--derivative element --- with $N(a)\le2\times10^9$ ($9.4\times10^8$ candidates, partners by full
factorisation); this is the expected order rather than a separate rigidity, a single fixed phase system
being heuristically $4^{\omega-1}$--fold rarer than the union of all systems. The contrast with
$\mathbb{Z}$ survives this accounting: over $\mathbb{Z}$ \emph{every} one of the $2^{\omega-1}$ sign
systems is barred by the theorem above, while over $\mathbb{Z}[i]$ the union of phase systems is
demonstrably populated from norm $1.7\times10^4$ on. The barrier of \#307 is thus localised precisely: it is the impossibility of phase
cancellation in the ordered field $\mathbb{R}$, where the only available phases $\pm1$ provably gain
nothing --- while the moment fourth roots of unity are admitted, the analogous existence problem closes
at seven primes. The cycles found come in conjugate pairs (supports not conjugation--stable); a
conjugation--\emph{stable} Gaussian cycle would descend to rational data through the norm form
$a^2+b^2$, i.e.\ precisely the Pythagorean layer of Proposition~\ref{prop:pairform} --- now a concrete
target rather than a speculation.
\end{remark}

\begin{remark}[The descended operator; the symmetric sector]\label{rem:descent}
Conjugation--stable supports admit an exact descent to $\mathbb{Z}$. Such a squarefree Gaussian integer
is an associate of an odd squarefree rational $m$ (split primes entering as whole conjugate pairs; the
$(1+i)$--sector is empty by a short parity argument), and a conjugation--equivariant unit assignment
makes the twisted derivative rational:
\[ D(m)\;=\;\sum_{\substack{p\mid m\\ p\equiv1\,(4)}} t_p\,\frac mp\;+\;\sum_{\substack{q\mid m\\
q\equiv3\,(4)}}\varepsilon_q\,\frac mq, \qquad t_p\in\{\pm2x_p,\pm2y_p\},\quad \varepsilon_q=\pm1, \]
where $p=x_p^2+y_p^2$ --- a Leibniz derivative on odd squarefree integers whose prime values have size
$2\sqrt p$ rather than $1$, and whose masses ${\sim}2/\sqrt p$ reach $2$ with four primes, so no
$59$--prime barrier exists for it. A two--cycle of $D$ is exactly a conjugation--stable Gaussian cycle
and would plant the phenomenon on rational soil. Parity constrains the hunt: $D(m)$ is odd iff $m$
carries an odd number of primes $\equiv3\pmod4$, so every member must contain such a prime. An
exhaustive search over members $m\le10^9$ (Gaussian norm $10^{18}$, $\omega\le8$, $1.4\times10^9$
closure tests) found no cycle, no fixed point $|D(m)|=m$, and no vanishing $D(m)=0$. We claim no
obstruction from this: equivariance prunes the available phases fourfold per split prime, and a phase
count makes emptiness at this depth the expected order. The descent target remains open over
$\mathbb{Z}[i]$ --- but not in the real--quadratic world, where the corresponding sector contains a fully
rational pair (Remark~\ref{rem:orderable}). Finally, the
descended coefficients are classical objects: at a split prime $p$ the four values
$\{\pm2x_p,\pm2y_p\}$ are exactly the Frobenius traces of the congruent--number curve $y^2=x^3-x$
(CM by $\mathbb{Z}[i]$) and of its twists --- $a_p\in\{\pm2x_p,\pm2y_p\}$ verified by point counting
for all $p<120$ --- so the conjugation--stable sector is the arithmetic derivative twisted by the
Hecke eigensystem of $\mathbb{Q}(i)$, and its mass statistics are governed unconditionally by Hecke's
equidistribution theorem: the first point of contact between \#307 and automorphic forms. Concretely:
over random squarefree $m$ the three--series conditions converge (the variance term is $\sum 2/p^2$), so
the descended mass $D(m)/m$ has an Erd\H{o}s--Wintner \emph{limiting distribution} whose prime components
are arcsine laws scaled by $1/\sqrt p$; the required input, equidistribution of the Hecke angles, is
visible exactly (Kolmogorov--Smirnov distance $0.0041$ from uniform over the $8{,}977$ split primes below
$2\times10^5$, well inside the $95\%$ band $0.0144$).
\end{remark}

\begin{remark}[The wall is the orderability of $\mathbb{Q}$]\label{rem:orderable}
The cycles over $\mathbb{Z}[i]$ permit a dissection of the barrier that no amount of theory over
$\mathbb{Z}$ could provide. The entire structural chain of this note --- rigidity, cross--matching, the
Pythagorean identity, the split discriminant --- transfers verbatim to twisted derivatives over
$\mathbb{Z}[i]$: Leibniz holds for disjoint supports, so a cycle $(a,b)$ has $D(N)=a^2+b^2$ for $N=ab$,
with $D(N)-2N=(a-b)^2$ and $2N-D(N)=\bigl(i(a-b)\bigr)^2$ (all verified exactly on the cycles found).
The last equation exposes the wall: over $\mathbb{Z}[i]$, where $-1=i^2$ is a square, \emph{both}
discriminant factors are squares automatically, and the plus/minus distinction of
Proposition~\ref{prop:split} --- which over $\mathbb{Z}$ carries the entire obstruction, via
$(a-b)^2\ge0\Rightarrow\sigma(N)\ge2$ --- has no content. The unique non--transferring link is the
sign--definiteness of squares: the barrier is, in the end, the \emph{formal reality} of $\mathbb{Q}$ in
the sense of Artin--Schreier. (The two Artin--Schreier theorems thus bracket the problem: their
characteristic--$2$ map builds the tower of Proposition~\ref{prop:ff-pyth}, their theory of orderable
fields names the wall.) This classifies where the mechanism can operate at all: it needs an ordering
making squares nonnegative (formal reality, hence a real embedding) \emph{and} bounded twisting phases
(finite unit group); by Dirichlet's unit theorem, unit rank $r_1+r_2-1=0$ with $r_1\ge1$ forces degree
one. \emph{$\mathbb{Q}$ is the only number field over which the mass barrier can operate}: over
imaginary quadratic fields squares lose sign--definiteness (over $\mathbb{Z}[i]$ even $-1$ is a square),
and over every other field the unit rank is positive and twisted masses are unbounded. This is a
classification of the wall's habitat, not an existence claim elsewhere --- but over $\mathbb{Z}[i]$
existence is no longer a claim: it is a list of sixteen. The ladder of imaginary quadratic rings
confirms the mechanism quantitatively. Over $\mathbb{Z}[\sqrt{-2}]$, whose only units are $\pm1$ ---
precisely the twisting freedom that provably gains nothing over $\mathbb{Z}$ --- sign--twisted cycles
nevertheless exist from norm $11{,}462$ on ($a=-90-41\sqrt{-2}$, $b=75-45\sqrt{-2}$, seven prime
classes): the complex embedding alone carries the phenomenon. Over the Eisenstein ring
$\mathbb{Z}[\omega]$, with six units, cycles appear at norm $2{,}964$ with \emph{six} prime classes
($a=-40+22\omega$, $b=56-21\omega$) --- the smallest derivative two--cycle exhibited in any ring ---
and the population at fixed height grows with the unit group: at norm $\le2\times10^7$ the censuses
read $14$, $56$, and $278$ twisted--cycle instances for unit groups of order $2$, $4$, $6$ --- a factor
${\sim}4$--$5$ per two extra units, as a phase count predicts --- while the strict canonical systems of
all three rings remain empty to norm $2\times10^9$ ($3\times10^9$ candidates in total, partners by full
factorisation). Two further objects appeared over $\mathbb{Z}[\sqrt{-2}]$: a composite with \emph{unit}
derivative ($a'\sim1$ --- impossible over $\mathbb{Z}$, where $n'=1$ characterises the primes), and an
\emph{anti--holomorphic} cycle: $a=-3163-1474\sqrt{-2}$ satisfies $a'=-\bar a$ exactly, under a single
conjugation--equivariant sign assignment closing both directions ($N(a)=N(b)=14{,}349{,}921$, five prime
classes each side). The equation $a'=u\bar a$ is vacuous over $\mathbb{Z}$, where it reads $n'=\pm n$;
it acquires solutions the moment conjugation is nontrivial. All cycles were verified exactly by
independent recomputation. Finally, the second pillar was tested in isolation: over
$\mathbb{Z}[\sqrt2]$ --- \emph{totally real}, so squares are sign--definite in both embeddings, but with
infinite unit group $\pm(1+\sqrt2)^{\mathbb{Z}}$ --- twisted cycles appear already at norm $882$: with
unit twists of height $\le2$, $a=21\sqrt2$ (classes $\sqrt2$, $-1+2\sqrt2$, $1+2\sqrt2$, $3$) cycles
exactly with $b=51+13\sqrt2$, of norm $2263$ and only \emph{two} prime classes --- against the
fifty--nine that $\mathbb{Z}$ demands --- and the composite $369-15\sqrt2$ has derivative exactly
$\varepsilon^6=99+70\sqrt2$, a unit. Both pillars of the classification are thus load--bearing, and each
is now confirmed by explicit counterexample on the soil where it alone fails: drop formal reality and
cycles appear (imaginary quadratic rings, norms $10^3$--$10^4$); keep formal reality but drop unit
finiteness and cycles appear (real quadratic, norm $882$, the smallest in any ring); over $\mathbb{Q}$,
where both hold, the barrier stands at $10^{112}$ by theorem. The bracketing of $\mathbb{Q}$ is complete
from both axes. The deep census sharpens it: at norm $\le2\times10^7$ the real--quadratic ring carries
$30{,}204$ twisted--cycle instances against $14$, $56$, $278$ for the imaginary rings with $2,4,6$ units
--- infinite unit rank dominates by two orders of magnitude --- and contains both the \emph{minimal}
architecture, $a=\sqrt2\,(5-\sqrt2)$ cycling with $b=7$ at norms $(46,49)$, two prime classes per member
(the rational prime $7$ being composite in the ring), the smallest derivative cycle found in any ring;
and a \emph{fully rational} pair:
\[ D(3707)=1547,\qquad D(1547)=3707\qquad(3707=11\cdot337,\ \ 1547=7\cdot13\cdot17), \]
exactly, with unit heights $\le2$ on both sides, verified independently: two ordinary integers in a
derivative two--cycle. The descent phenomenon that is empty over $\mathbb{Z}[i]$ to norm $10^{18}$ is
realised over $\mathbb{Z}[\sqrt2]$ at four digits. The canonical system of $\mathbb{Z}[\sqrt2]$ remains
cycle--free to norm $2\times10^9$ ($7.7\times10^8$ candidates), now with seven unit--derivative
composites.
\end{remark}

The cycle phenomenon over rings with units is governed by the theory of unit equations, which supplies
its rigidity theorem.

\begin{proposition}[$S$--unit rigidity of twisted cycles]\label{prop:sunit}
Let $\mathcal{O}$ be the ring of integers of a number field and $S$ a finite set of primes of
$\mathcal{O}$. If $(a,b)$ is a twisted two--cycle with $\operatorname{supp}(ab)\subseteq S$, then
dividing $D(a)=b$ by $b$ exhibits a solution of
\[ x_1+\dots+x_{\omega(a)}=1,\qquad x_i=u_i\,\frac{a/\pi_i}{b}, \]
in the $S$--unit group of the field, and likewise for $b$. By the theorem of
Evertse--Schlickewei--Schmidt \cite{ESS2002} (cf.\ \cite{Evertse1984}), such an equation has only
finitely many nondegenerate solutions; hence \emph{every fixed support carries at most finitely many
twisted cycles free of vanishing subsums}, and the species can be infinite only through growing
supports --- precisely the behaviour of every census above. (Vanishing subsums correspond to partial
zero--derivative patterns; none occurred in any search. The relation was verified exactly on the
$\mathbb{Z}[\sqrt2]$ specimen: the four $S$--unit terms of $a=21\sqrt2$ sum to $1$.)
\end{proposition}

\begin{remark}[Galois folding, and the thinness of symmetric sectors]\label{rem:folding}
Over a quadratic field with automorphism $\sigma$, a $\sigma$--equivariant choice of unit values
($u_{\pi^\sigma}=u_\pi^{\,\sigma}$) makes $D$ commute with conjugation, so the \emph{single} equation
$D(a)=a^\sigma$ already implies that $(a,a^\sigma)$ is a twisted two--cycle (checked symbolically and on
$1{,}988$ random instances). In the two--prime case the partner is forced:
$\rho=-\bigl(\bar u v\,\pi+\bar v\,(\pi^\sigma)^2\bigr)/\bigl(N(u)-N(\pi)\bigr)$, with integrality a
congruence and the entire residue of the problem the primality of $|N(\rho)|$ along Pell--exponentially
growing sequences --- Lehmer territory rather than Bunyakovsky. Whether a folded sector is starved or
rich is decided by its degrees of freedom. When the fold leaves no slack the sector is thin: the folded
\emph{two}--prime family over $\mathbb{Z}[\sqrt2]$ is empty in the tested box ($36{,}000$ parameter
triples, none even integral), the conjugation--stable sector over $\mathbb{Z}[i]$ is empty to Gaussian
norm $10^{18}$, and the anti--holomorphic cycle over $\mathbb{Z}[\sqrt{-2}]$ is unique in its census.
But with \emph{three} primes the folded equation
$u_1\pi_2\pi_3+u_2\pi_1\pi_3+u_3\pi_1\pi_2=(\pi_1\pi_2\pi_3)^\sigma$ carries three unit unknowns against
two real equations, and this underdetermined system is \emph{populated}: an exhaustive box ($p_i\le137$,
unit heights $\le4$) contains twenty--four solutions over eleven prime triples, the smallest being
\[ a=(3+\sqrt2)(5-2\sqrt2)(5-\sqrt2)=57-16\sqrt2,\qquad D(a)=a^\sigma=57+16\sqrt2, \]
with unit values $\varepsilon^2$, $\varepsilon$, $-\varepsilon^{-1}$: \emph{the derivative conjugates the
element}, at norm $2737$ on both sides, the equivariant return being automatic (verified exactly).
Existence of twisted cycles over the four rings tested is a theorem by exhibition; for \emph{every} ring
with infinite unit group it is now reduced to the abundance of lattice points on these underdetermined
unit systems --- measured plentiful already at $k=3$ --- a geometry--of--numbers question rather than a
primality one.
\end{remark}

We codify what is theorem and what remains open in this circle.

\begin{theorem}[Existence of twisted derivative two--cycles]\label{thm:existence}
Leibniz derivatives with unit prime--values admit genuine two--cycles over $\mathbb{Z}[i]$,
$\mathbb{Z}[\sqrt{-2}]$, $\mathbb{Z}[\omega]$, and $\mathbb{Z}[\sqrt2]$. The proof is by exhibited
certificate, each verified exactly by independent recomputation: the smallest instances are, in order of
norm, $\bigl(\sqrt2(5-\sqrt2),\,7\bigr)$ over $\mathbb{Z}[\sqrt2]$ at norms $(46,49)$; the Eisenstein
cycle at norm $2{,}964$; the sign--twisted cycle over $\mathbb{Z}[\sqrt{-2}]$ at norm $11{,}462$; and the
cycle over $\mathbb{Z}[i]$ at norm $16{,}642$ --- together with the rational pair
$D(3707)=1547$, $D(1547)=3707$ over $\mathbb{Z}[\sqrt2]$.
\end{theorem}

\begin{proposition}[Geometry of the solution set]\label{prop:geometry}
Over $\mathbb{Z}[\sqrt2]$, consider the Galois--folded system
$u_1\pi_2\pi_3+u_2\pi_1\pi_3+u_3\pi_1\pi_2=(\pi_1\pi_2\pi_3)^\sigma$ of
Remark~\textup{\ref{rem:folding}}. \textup{(i)} For \emph{every} triple of split primes the system is
solvable with the $u_i$ on the real unit hyperbola $\{|st|=1\}$: solving the second embedding for $u_2$
and letting the third coordinate tend to zero, the first--embedding residual sweeps $\pm\infty$ while the
remaining terms stay bounded, so a root exists by continuity; the real solution set is a nonempty curve
\textup(verified numerically on $200$ random triples, all solvable\textup). The volume of the solution
variety never vanishes. \textup{(ii)} By Proposition~\textup{\ref{prop:sunit}} each fixed support carries
finitely many integral solutions, so infinitude can only come from growing supports. \textup{(iii)} The
measured integral hit rates are $11$ of $455$ prime triples in the smallest box, and $30{,}204$ cycles of
norm $\le2\times10^7$ at twist height $\le2$. \textup{(iv)} What remains open is precisely whether the
discrete unit lattice meets the real curve in integer points for infinitely many supports ---
inhomogeneous Diophantine approximation on exponentially parametrised curves. Two shortcuts are closed:
bounded--unit polynomial families are obstructed (the rational component of the return forces unit
coefficients to track growing prime coordinates, which bounded units cannot), and even the cleanest
forward family --- $D\bigl(\sqrt2(x-\sqrt2)\bigr)=x+2$ \emph{identically in} $x$, supplying a candidate
cycle whenever $x^2-2$ and $x+2$ are prime --- has sporadic returns: the closing two--term unit equation,
effectively finite for each $x$, is solvable at $x=5$ and provably insolvable at $x=15$ and $x=21$.
Those failures are, however, exactly the rigid case: a \emph{prime} return $x\pm2$ leaves the closing
system no slack. The identity has four unit branches $D\bigl(\sqrt2(x\pm\sqrt2)\bigr)=x\pm2$ (independent
signs, twists $1$ and $\pm\varepsilon^{\pm1}$), and when $x\pm2$ is composite squarefree the closing
system gains one unit unknown per prime class against the same two equations --- and does close:
$x=107,121,387,1693$ yield cycles with rational partners $b=105,119,385,1695$, each verified exactly.
Infinitude along this family is thereby reduced to Bunyakovsky for $x^2-2$ together with recurring
composite--return closure: the softest form the gap takes anywhere in this paper. Even the primality
half of that reduction can be removed. By the half--dimensional sieve
\cite{Iwaniec1978,LemkeOliver2012} (applied to $(2y+1)^2-2=4y^2+4y-1$ to keep $x$ odd), $x^2-2$ is a
$P_2$ infinitely often; every prime factor of $x^2-2$ has $2$ as a quadratic residue and therefore
splits in $\mathbb{Z}[\sqrt2]$, so $\sqrt2(x-\sqrt2)$ is infinitely often squarefree and all--split
with $\omega\in\{2,3\}$. The candidate supports thus exist \emph{unconditionally}; all that remains
open on this route is whether the unit gates open along that supply infinitely often.
Exactness migrates; it does not disappear.
\end{proposition}

\begin{proposition}[No vanishing derivatives, no transport, and the price of exactness]\label{prop:novanish}
\textup{(i)} No squarefree $z$ with $\omega(z)\ge1$ has $D(z)=0$, for any twisted derivative over any of
the rings considered (indeed over any UFD): modulo a prime $\pi\mid z$ every cofactor term but one
vanishes, leaving $D(z)\equiv u_\pi\,z/\pi\not\equiv0\pmod\pi$. The zero--derivative counters of every
census above are therefore theorems, not observations. \textup{(ii)} Consequently cycles cannot be
transported multiplicatively: for squarefree $c$ coprime to $ab$, Leibniz gives $D(ca)=D(c)\,a+c\,D(a)$,
so $(ca,cb)$ is a cycle exactly when $D(c)=0$ --- impossible. Cycles do not breed; each one requires
fresh exactness. \textup{(iii)} The two cycle equations can be bought by construction, at a conserved
price. Buying \emph{neither} leaves the generic stratum: two exact unit conditions ($30{,}204$
specimens). Buying \emph{one} is free in two dual ways --- forward, $\pi\coloneqq q-\varepsilon^m\sqrt2$
makes $D(\sqrt2\,\pi)=q$ an identity (the four--branch family above is the case $m=\pm1$); return,
$\pi\coloneqq(\rho+\varepsilon^{-1}\rho^\sigma)/\sqrt2=(c+2d)-d\sqrt2$ for $\rho=c+d\sqrt2$ closes
$D(b)=\sqrt2\,\pi$ identically for every associate $b$ of $\rho\rho^\sigma$ --- and the residue is one
primality, $|N(\pi)|$, plus one exact unit gate, measured at $2$ of $50$ pairs with both primalities in
the smallest box. The two hits re--derive, with no search, precisely the two census minimal--shape cycles
with prime rational partner: norms $(46,49)$ and $(206,49)$. Buying \emph{both} collapses the parameters
to unit--indexed quadratics in $q$ with a square--discriminant condition; the resulting $q=7$ line
carries the known cycles at $m=1$ and $m=-2$ (discriminants $25$ and $81$) together with one Galois
image, and every other prime--norm rung out to $|m|\le14$ fails its return, verified exactly. Each
purchased identity is repaid as a sparser arithmetic condition elsewhere: the hardness of exactness is
conserved, never destroyed.
\end{proposition}

\begin{remark}[The two faces of the gap: unit rank dichotomises the obstruction]\label{rem:dichotomy}
The invariant that governs abundance in Remark~\ref{rem:orderable} --- the unit rank
$r=r_1+r_2-1$ of Dirichlet --- also governs the \emph{type} of the residual open problem, and the two
types are mutually exclusive. When $r=0$ (the rings $\mathbb{Z}$, $\mathbb{Z}[i]$,
$\mathbb{Z}[\sqrt{-2}]$, $\mathbb{Z}[\omega]$) the unit group is finite, so a given pair of supports
admits only $|\mathcal{O}^\times|$ twist assignments and hence finitely many candidate cycles: the
solution locus is a set of \emph{isolated points}, zero--dimensional, and infinitude can arise only from
infinitely many \emph{distinct} supports independently passing a closing test --- a
Bateman--Horn/Schinzel--type recurrence over prime tuples, with no continuous family to deform along.
The constituent primalities are individually classical (each support prime is a value of a binary
norm--form, so by Fermat--Dirichlet --- and, for $\mathbb{Z}[i]$, by the Green--Sawhney theorem with
\emph{both} coordinates prime --- there is no shortage of supports); it is their \emph{simultaneous}
closure that is not classical. This was tested directly: the minimal shape $a=(1+i)\pi$ over
$\mathbb{Z}[i]$ yields \emph{no} cycle for any Gaussian prime $\pi$ of norm $\le4000$ under any of the
sixteen twists, and the sixteen genuine cycles are irregular mixtures of inert and split classes
($\omega=3,4,5$) with no one--parameter family among them --- the signature of isolated points, not a
curve. When $r\ge1$ (every real quadratic ring, e.g.\ $\mathbb{Z}[\sqrt2]$),
Proposition~\ref{prop:geometry} makes the real solution locus a positive--dimensional curve in each
support, but the unit lattice is an infinite rank--$r$ discrete set, and infinitude requires it to meet
that curve in \emph{integer} points infinitely often --- inhomogeneous Diophantine approximation on an
exponentially parametrised curve. No unit rank delivers both a classical gate and a continuous family:
rank zero gives the former and denies the latter, positive rank the reverse. This is why no single ring
lets the infinitude be \emph{finished} by present methods --- a structural account of the obstruction's
universality, not a barrier theorem. It also locates $\mathbb{Z}$ exactly: the unique ring carrying
\emph{both} the rank--zero sporadicity \emph{and}, by formal reality, the mass barrier --- the worst
cell of the trichotomy of Corollary~\ref{cor:ff307}, which is precisely why \#307 is the hard
representative of its family.
\end{remark}

\section{A Th\=abit--type rule and its honest limits}\label{sec:frame}

The bridge recasts the search for a solution as the search for a derivative two--cycle, and the
derivative's Leibniz structure makes the cycle equations \emph{linear} in the ``closing'' primes.
This yields a constructive calculus analogous to Th\=abit ibn Qurra's rule for amicable pairs and to
Euler's divisor method.

\begin{theorem}[Frame Rule]\label{thm:frame}
Let $M,N$ be coprime squarefree integers with derivatives $M',N'$, and set $\Delta\coloneqq MN-M'N'$.
Suppose
\[
p=\frac{MN'+N^2}{\Delta},\qquad q=\frac{M'N+M^2}{\Delta}
\]
are integers, both prime, each coprime to $MN$, with $p\ne q$. Then $a\coloneqq Mp$, $b\coloneqq Nq$
form a derivative two--cycle $a'=b$, $b'=a$; equivalently $(P,Q)$ with $P$ the prime set of $Mp$ and
$Q$ that of $Nq$ solves Erd\H{o}s \#307.
\end{theorem}

\begin{proof}
With $p$ prime and $p\nmid M$, $a=Mp$ is squarefree and $a'=M'p+M$; similarly $b'=N'q+N$. The cycle
conditions $a'=b$, $b'=a$ read
\[
M'p+M=Nq,\qquad N'q+N=Mp,
\]
a linear system in $(p,q)$ with determinant $-(MN-M'N')=-\Delta$. Solving gives exactly
$p=(MN'+N^2)/\Delta$ and $q=(M'N+M^2)/\Delta$; substituting back verifies both equations
identically. (Both identities were checked symbolically.)
\end{proof}

A two--primes--per--side variant (``Rule~2'') is obtained identically: with $a=Mp_1p_2$, $b=Nq_1q_2$
and elementary symmetric functions $e_1,e_2,f_1,f_2$, the cycle equations are linear in the $f_i$,
\[
f_2=\frac{Me_1+M'e_2}{N},\qquad f_1=\frac{MNe_2-MN'e_1-M'N'e_2}{N^2},
\]
and a solution is realised whenever $f_1^2-4f_2$ is a perfect square with prime roots.

\paragraph{An Euler-style divisor trick.}
Fix $N$ and let $M=M_0r$ with $r$ a new prime, so $M'=M_0'r+M_0$. Put
$\Delta_0\coloneqq M_0N-M_0'N'$.

\begin{proposition}[Divisor trick]\label{prop:divisor}
With the above notation and $p_{\mathrm{num}}\coloneqq MN'+N^2$:
\begin{enumerate}[label=\textup{(\alph*)},leftmargin=2.2em]
\item $\Delta(r)\coloneqq MN-M'N'=r\,\Delta_0-M_0N'$ \emph{(linear in $r$)};
\item $\Delta_0\,p_{\mathrm{num}}-M_0N'\,\Delta(r)=K$ where $K\coloneqq \Delta_0N^2+(M_0N')^2$ is
\emph{constant in $r$};
\item consequently $\Delta(r)\mid p_{\mathrm{num}}\Rightarrow \Delta(r)\mid K$, and the converse
holds provided $\gcd(\Delta(r),\Delta_0)=1$ \textup{(}equivalently $\gcd(M_0N',\Delta_0)=1$, since
$\Delta(r)\equiv-M_0N'\bmod\Delta_0$\textup{)};
\item each positive divisor $d\mid K$ with $d\equiv-M_0N'\ (\mathrm{mod}\ \Delta_0)$ yields a
candidate $r=(d+M_0N')/\Delta_0$ and recovers $\Delta(r)=d$; this enumeration is complete (it
captures every genuine solution). The resulting $p=p_{\mathrm{num}}/d$ is integral when
$\gcd(M_0N',\Delta_0)=1$, and otherwise must be checked.
\end{enumerate}
\end{proposition}

\begin{proof}
(a) and (b) are direct expansions (verified symbolically; $\partial K/\partial r=0$). (c): from
(b), $\Delta(r)\mid p_{\mathrm{num}}\Rightarrow\Delta(r)\mid\Delta_0p_{\mathrm{num}}
=M_0N'\Delta(r)+K\Rightarrow\Delta(r)\mid K$. For the converse, $\Delta(r)\mid K$ gives
$\Delta(r)\mid\Delta_0p_{\mathrm{num}}$, whence $\Delta(r)\mid p_{\mathrm{num}}$ iff
$\gcd(\Delta(r),\Delta_0)=1$. (d) inverts (a). Excluded degeneracies: $\Delta_0=0$ ($r$ undefined)
and $K\le0$.
\end{proof}

Thus, exactly as in Euler's amicable--pair machinery, one factors the single number $K$ once and
reads off candidates from its divisors lying in a fixed residue class modulo $\Delta_0$; only the
primality of $r,p,q$ then remains. Two caveats deserve emphasis:
the converse in (c) and the automatic integrality of $p$ require $\gcd(M_0N',\Delta_0)=1$ (it fails,
e.g., for $M_0=2$, $N=19\cdot47$, $r=5$, where $\Delta(r)\mid K$ but $\Delta(r)\nmid
p_{\mathrm{num}}$); and the count of admissible classes is governed by $\Delta_0$, not merely by
$\tau(K)$.

When the closing side carries a \emph{single} new prime, the rule collapses to a normal form that
separates the two layers of the problem.

\begin{proposition}[One--prime normal form; the exactness stratum]\label{prop:oneprime}
\textup{(a)} Erd\H{o}s \#307 has a solution if and only if there exist squarefree $M,N$ with
\[ MN-M'N'=M^2, \]
such that $p\coloneqq N'/M$ is a prime not dividing $M$. Indeed the equation forces $M\mid M'N'$, hence
$M\mid N'$ by $\gcd(M,M')=1$, and $N=M+M'p$; then $a=Mp$, $b=N$ satisfy $a'=M'p+M=N$ and $b'=N'=Mp=a$
(coprimality and $a\ne b$ are automatic). Conversely every solution $(a,b)$ and every prime $p\mid a$
yield such a pair with $M=a/p$, so each solution carries $\omega(a)$ witnesses. The problem thus splits
into an \emph{exactness} layer (the Diophantine equation) and a single \emph{primality} layer (one prime
value of the determined quantity $N'/M$).
\textup{(b)} \emph{The exactness layer alone carries the full barrier.} Call $(M,N)$ a \emph{relaxed
witness} if $MN-M'N'=M^2$ with $p=N'/M$ composite. The identity
$\bigl(\sigma(M)+\tfrac1p\bigr)\sigma(N)=\frac{N}{Mp}\cdot\frac{Mp}{N}=1$ holds for \emph{every} witness,
so by strict AM--GM ($\sigma(N)=1$ being impossible) $\sigma(M)+\sigma(N)>2-\tfrac1p$. Rigidity forces
three structural facts. The supports of $M$ and $N$ are \emph{automatically} disjoint: a common prime
$q$ would divide $N'=Mp$ through $q\mid M$, against $\gcd(N,N')=1$. Likewise $\gcd(p,MN)=1$: a prime
shared by $p$ and $M$ propagates to $N=M+M'p$ and lands in the previous case, while one shared by $p$
and $N$ divides $N'$. And if $p$ is even, both members are odd: $2\mid N'=Mp$ forces $N$ odd by the same
lemma, while $M$ even would make $N=M+M'p$ even. The union of supports therefore avoids every prime
divisor of $p$, and for even $p$ avoids $2$; exact greedy computation then gives, for \emph{every}
composite $p$,
\[ \omega(MN)\ \ge\ 59, \qquad MN\ >\ 8.77\times10^{112} \]
--- the cycle barrier itself --- with strictly stronger floors at every finite $p$: $\omega\ge194$,
$MN>10^{497}$ at $p=9$; $\omega\ge230$, $MN>10^{609}$ at $p=4$; the even--$p$ floors rising towards the
all--odd bound $10^{5040}$; and the infimum $8.77\times10^{112}$ approached only as $p\to\infty$ through
composites free of prime factors $\le277$. Deleting primality thus lowers the barrier by
\emph{nothing}: the $59$--prime, $10^{112}$ wall is carried entirely by the exactness equation, and the
primality of $N'/M$ is a separate demand on top of it. (No witness, prime or composite, exists with
$M\le120$, $p\le2000$, exhaustively.)
\end{proposition}

\paragraph{The decisive obstruction: a circular objective.}
The heuristic yield of the divisor trick from one frame is
\[
\mathbb{E}[\text{cycles}] \ \approx\ A\cdot\frac{1}{\ln r\,\ln p\,\ln q},
\qquad A=\#\{d\mid K:\ d\equiv -M_0N'\ (\mathrm{mod}\ \Delta_0)\}\ \approx\ \frac{\tau(K)}{\Delta_0},
\]
the last step being an equidistribution heuristic. To make $\mathbb{E}\ge1$ at the barrier scale
($1/\ln^3\approx 5\times10^{-7}$) one needs $A\gtrsim 2\times10^6$ admissible divisors per frame,
hence $\Delta_0=O(1)$ with $\tau(K)\sim10^6$. The natural objective is therefore ``engineer frames
with $\Delta_0$ small.'' But here is the obstruction, an exact identity:
\[
\boxed{\ \Delta_0 \;=\; M_0N\bigl(1-s_{M_0}\,s_N\bigr),\qquad
s_{M_0}=\!\!\sum_{p\mid M_0}\!\frac1p,\ \ s_N=\sum_{p\mid N}\frac1p.\ }
\]
So $\Delta_0$ is small \emph{relative to} $M_0N$ exactly when $s_{M_0}s_N\approx1$ --- which is the
Erd\H{o}s~\#307 condition for the frames themselves --- and $\Delta_0=O(1)$ in absolute terms at
scale requires $s_{M_0}s_N=1-O(1/M_0N)$, i.e.\ a reciprocal--product within $O(1/M_0N)$ of $1$,
\emph{strictly harder} than \#307. The frame--engineering objective is thus circular: the Frame Rule
is a faithful \emph{parametrisation} of solutions, not a reduction in difficulty.

\begin{remark}[Honest status of the constructive program]
The Frame Rule (Theorem~\ref{thm:frame}) and divisor trick (Proposition~\ref{prop:divisor}) are
exact theorems; they make the existence of a solution \emph{falsifiable in principle by finite
computation per frame}, and they place \#307 in the same constructive lineage as Th\=abit's and
Euler's rules. They do not, as they stand, lower the difficulty below \#307 itself: by the boxed
identity the bottleneck (small $\Delta_0$) is the original problem in disguise, and a complete
small--frame sweep ($2.2\times10^5$ frames) produces zero cycles with total expected count $<1$, as
the barrier demands. We regard the calculus as an exact parametrisation of solutions whose central
obstruction we have now identified precisely --- not a pathway that lowers the difficulty, and not
evidence that a solution is within reach.
\end{remark}

\section{A breeder form and the feasibility of certificate search}\label{sec:breeder}

Allowing \emph{two} free primes on one side turns the divisor trick into a cleaner bilinear form ---
Euler's amicable--pair method, in the shape used by te~Riele's breeding algorithms
\cite{teRiele1984} --- and lets us state precisely the strongest honest conclusion: a
\emph{conditional} reduction, which nonetheless does not become a feasible finite search.

\begin{proposition}[Bilinear breeder form]\label{prop:bde}
Fix coprime squarefree $M_0,N$ and seek a two--cycle $a=M_0rp$, $b=Nq$ with $r,p,q$ distinct primes
coprime to $M_0N$. Put $G:=NM_0-N'M_0'\ (=\Delta_0)$, $c:=M_0N'$, and $K:=GN^2+c^2$. Then the cycle
equations $a'=b$, $b'=a$ are equivalent to
\[
(Gr-c)(Gp-c)=K, \qquad q=\frac{M_0rp-N}{N'} .
\]
\end{proposition}
\begin{proof}
Expanding, $(Gr-c)(Gp-c)-K=-N'G\,(a'-b)$ (verified symbolically), so the bilinear equation is
equivalent to $a'=b$; the second cycle equation $b'=a$ then determines $q$ as stated.
\end{proof}

Each factorisation $K=d\cdot(K/d)$ with $d\equiv-c\ (\mathrm{mod}\ G)$ proposes a candidate
$(r,p)=\bigl((d+c)/G,(K/d+c)/G\bigr)$; a solution is any such candidate whose $r,p,q$ are all prime.
A dual, partition--side form is occasionally useful:

\begin{proposition}[Subset form]\label{prop:subset}
For a support $U$ with $D=\prod_{p\in U}p$ and $\Sigma=\sum_{p\in U}D/p$, a partition $U=P\sqcup Q$
solves \#307 iff $S=P$ satisfies $A(\Sigma-A)=D^2$, where $A:=\sum_{p\in S}D/p$. Moreover
$A(\Sigma-A)-D^2=-(1-s_Ps_Q)\,D^2$.
\end{proposition}

(Both identities verified symbolically and on $2\times10^4$ random supports.) The last identity says
the ``analytic'' residual $1-s_Ps_Q$ and the ``arithmetic'' residual $A(\Sigma-A)-D^2$ are one
equation up to the factor $D^2$: the two obvious steering handles carry no independent information.

\paragraph{A conditional reduction.} Several ingredients of a construction are unconditional. Greedy
Egyptian--fraction choice gives disjoint $P,Q$ with $0<1-s_Ps_Q=O(1/T)$ (solvability to any
precision); a prime $r$ nearest $M_0N'/\Delta_0$ gives $G\le\Delta_0^{0.475}(M_0N')^{0.525}$
unconditionally (Baker--Harman--Pintz \cite{BHP2001}); and each small prime $\ell$ can be forced to
divide $K$ by one congruence on a free prime slot of $N$ (Dirichlet then supplies the prime).
Granting these, \#307 reduces to a single prime--density statement,
\[
(\mathrm{H})\qquad\text{among the admissible candidates so produced, some triple $(r,p,q)$ is prime,}
\]
and $(\mathrm{H})\Rightarrow$ \#307 is \textsc{yes}. This $(\mathrm{H})$ is a
Hardy--Littlewood/Bateman--Horn--type assertion for an explicit, computable, \emph{non--polynomial}
family; we know of no named conjecture implying it. That is not an accident of our construction:

\begin{proposition}[No polynomial families]\label{prop:nopoly}
Let $p_1(t),\dots,p_k(t)$, $q_1(t),\dots,q_l(t)$ be integer polynomials, each taking prime values for
infinitely many integers $t$ (a fixed--shape Schinzel--type family), such that
$\bigl(\sum_i1/p_i(t)\bigr)\bigl(\sum_j1/q_j(t)\bigr)=1$ for infinitely many $t$. Then every $p_i$ and
$q_j$ is constant: the family is a single fixed solution. Consequently no polynomial parametrisation can
reduce \#307 to Bunyakovsky, Schinzel's Hypothesis~H, or Bateman--Horn.
\end{proposition}
\begin{proof}
Holding at infinitely many $t$, the relation holds identically as rational functions. A nonconstant
$p_i$ taking infinitely many prime (hence positive) values has positive leading coefficient on an
infinite branch, so $p_i(t)\to+\infty$ there; split each side into constants and nonconstant members,
$\sum_i1/p_i=A_P+\varepsilon_P(t)$ with $\varepsilon_P(t)\to0^+$. Letting $t\to\infty$ in the identity
gives $A_PA_Q=1$ (in particular both sides contain constants), and subtracting,
$A_P\varepsilon_Q(t)+A_Q\varepsilon_P(t)+\varepsilon_P(t)\varepsilon_Q(t)=0$ for all large $t$. Every
term on the left is nonnegative, and strictly positive as soon as one nonconstant member exists; hence
there are none.
\end{proof}

\noindent This upgrades the earlier computational observation (that the breeder family is not
polynomially parametrised) to a theorem for \emph{all} fixed--shape families: the prime--density
hypothesis behind \#307 is irreducibly non--polynomial, which is why no standard conjecture applies.

\paragraph{Why $(\mathrm{H})$ is not a finite search.} It is tempting to read $(\mathrm{H})$ as
``run a certificate hunt that is expected to succeed.'' It is not, for a quantitative reason. The
expected number of certificates from one frame is $\mathbb{E}\approx
\bigl(\tau(K)/G\bigr)\cdot(\ln r\,\ln p\,\ln q)^{-1}$, and \emph{both} factors are pinned by the
geometry.
\begin{itemize}[leftmargin=1.4em]
\item \emph{Demand.} $G=\Delta_0=M_0N(1-s_{M_0}s_N)$ exactly, so $G$ is small only in the circular
regime $s_{M_0}s_N\approx1$; over $2\times10^4$ random coprime squarefree frames the ratio
$G/(M_0N)$ has median $1$ and best value $\approx10^{-0.23}$, and the honest Baker--Harman--Pintz
reduction still leaves $\log_{10}G\approx13$ at the barrier.
\item \emph{Supply.} $\tau(K)$ is bounded by the \emph{size} of $K=GN^2+c^2\ (\ge c^2)$, which is an
\emph{output} of the frame, not a free dial. The maximal order of the divisor function
\cite[Thm.~317]{HardyWright} caps $\log_2\tau(K)$ at $\approx36$ for $K\le10^{54}$ (the barrier
geometry) and at $\approx112$ even for $K\le10^{240}$; reaching $\tau(K)=2^{120}$ needs
$K\gtrsim10^{263}$, i.e.\ frame scale far beyond the barrier.
\end{itemize}
Enlarging a frame grows $K\propto N^2$ but gains divisors only at rate $\log2/\log\log K<1$ per
log--unit of $K$, while it grows $G$ at rate $1$ per log--unit of $N$; so $\tau(K)/G$ cannot be
driven up. Optimising $\mathbb{E}$ over all $(M_0,N)$ splits with geometry--consistent $K$ and the
maximal--order $\tau$, and granting the most favourable demand and singular--series constants, yields
$\log_{10}\mathbb{E}\le-4$ at the barrier and \emph{decreasing} with scale; empirically $385$
factorable frames realise a median of $0$ and a maximum of $1$ admissible divisor, against the
$\sim10^{5}$ a single productive frame would need. A self--consistent accounting --- capping the
number $w$ of forced prime factors of $K$ by the frame's own residue freedom (counted generously) ---
sharpens this to $\log_{10}\mathbb{E}\le-8$ at $B=56$, falling to $\approx-260$ at $B=3000$: the
deficit is \emph{scale--invariant}, so no choice of scale, slot count, or budget split reaches
expectation $1$. (The construction is even self--similar: engineering $K$'s algebraic --- e.g.\
Gaussian--integer --- factorisation to supply divisors reduces to prescribing $N$ and $N'$
simultaneously, an arithmetic--antiderivative inverse problem of the same difficulty class as \#307.)
The breeder thus makes \#307 conditional on $(\mathrm{H})$ but does \emph{not} render it a feasible
finite computation: the supply--demand conflict through $K=GN^2+c^2$ is exactly the circularity of
Section~\ref{sec:frame} in quantitative form. There is also a structural reason the analogy with
amicable pairs was bound to fall short: te~Riele's breeding amplifies a \emph{known seed} pair into
new ones, whereas the barrier (Theorem~\ref{thm:barrier}) forbids any seed below $2.09\times10^{56}$
--- the method imports a precondition the problem cannot supply.

\section{Computations}\label{sec:computations}

All verdicts use exact integer/rational arithmetic; floating point is used only for pre--screening.
\begin{itemize}[leftmargin=1.4em]
\item \emph{Exhaustive non--existence below $10^8$.} Enumerating all prime sets $P$ with
$\sum1/p>1$ and $\prod P\le10^8$ (forcing $Q$ as the prime support of $N_P$ and testing
$N_Q=D_P$): $1{,}075{,}419$ candidate sets, no solution. Subsumed by Theorem~\ref{thm:barrier} but an
independent sanity check.
\item \emph{Derivative sieve.} No solution of $n''=n$ with $n'\ne n$ and both members $\le
2\times10^7$ (smallest--prime--factor sieve of the derivative).
\item \emph{Function field.} The arithmetic derivative on $\mathbb{F}_q[t]$ has no two--cycles and no
nonzero fixed points, verified exhaustively over $\mathbb{F}_2$ (degree $\le7$), $\mathbb{F}_3$
($\le5$), $\mathbb{F}_5$ ($\le4$), confirming Theorem~\ref{thm:ff}.
\item \emph{Density.} The sieve gives density $0.042029$ to $2\times10^7$ ($0.017623$ among
squarefree); the independent Erd\H{o}s--Wintner convolution over primes $\le10^5$ gives
$\delta=0.04202$, agreeing to four places (Proposition~\ref{prop:density}).
\item \emph{Conditional mass.} Of $43{,}087$ qualifying squarefree $a$ sampled, $b=a'$ is squarefree
$95.0\%$ of the time; on the resulting $40{,}920$ squarefree $b$ the mean reciprocal mass is $0.047$
vs.\ the required $1/s\approx0.934$, with empirical tail probability $\Pr[\sum_{q\mid b}1/q\ge1/s]=
0/40{,}920$.
\item \emph{Frame/ledger sweep.} Over $2.2\times10^5$ factorable small frames: the divisor--trick
identities hold on every instance; the smallest $\Delta_0$ observed is $22$; the equidistribution
$A\approx\tau(K)/\Delta_0$ holds only in aggregate ($\sum A=127$ vs.\ $\sum\tau(K)/\Delta_0=39.8$)
and fails per frame ($A=0$ for $99.95\%$ of frames); total expected cycles $<1$; zero cycles, as the
barrier requires.
\item \emph{$\mu$--Sondow line census.} Enumerating the lines $n'=n\pm d$ (squarefree members coprime to
$d$) finds no adjacent opposite--sign pair at additive distance $d$: exhaustively to $2\times10^7$ for
$|d|\le200$, and along the small--$|\delta|$ genealogy tree out to members of more than a thousand digits.
Every $\delta=\pm1$ member encountered is a known Giuga or primary pseudoperfect number; the closest
opposite--sign approach recorded is $|y-(x+d)|=11$ (at $d=1$, $30$ vs.\ $42$, and $d=37$, $210$ vs.\ $258$),
attained only among small members --- the two trees do not come close at height.
\end{itemize}
Code and logs accompany this note.

\section{Formalisation in Lean~4}\label{sec:lean}

The rigidity and barrier results are formalised against mathlib (Lean~4, toolchain
\texttt{v4.30.0}). The following are machine--checked and \texttt{sorry}--free, depending only on the
standard axioms \texttt{propext}, \texttt{Classical.choice}, \texttt{Quot.sound}:
\begin{itemize}[leftmargin=1.4em]
\item Lemma~\ref{thm:rigidity} (\texttt{rigidity\_coprime}) and Theorem~\ref{thm:structure}
(\texttt{solution\_structure});
\item Lemma~\ref{lem:amgm} (\texttt{reciprocal\_sum\_gt\_two}); the exact thresholds
$\sum_{\text{first }58}1/p<2<\sum_{\text{first }59}1/p$ grounding $|U|\ge59$
(\texttt{recipSum58\_lt\_two}, \texttt{recipSum59\_gt\_two}); the numeric core
$\Pi_{59}^2\ge(4\times10^{112})N_{59}$ (\texttt{barrier\_numeric}); and the algebraic core of the
barrier (\texttt{barrier}), giving $D_P^2\ge4\times10^{112}$.
\end{itemize}
The single non--elementary ingredient---that the $k$ smallest primes minimise $\prod$ and maximise
$\sum1/p$ among $k$--element prime sets---is isolated as an explicit hypothesis of \texttt{barrier}
(the ratio bound $R/T(U)\ge\Pi_{59}/T_{59}$), so the dependency is transparent rather than hidden;
its standalone proof via \texttt{Nat.nth\,Nat.Prime} is routine and left as future work.

\section{Status and questions}\label{sec:status}

Erd\H{o}s~\#307 remains \textbf{open}. We have shown it is equivalent to the existence of a non--trivial
two--cycle of the arithmetic derivative; established rigidity and a barrier
(Rosen's $|P\cup Q|\ge59$, together with our $\min(\prod P,\prod Q)\ge2.09\times10^{56}$) that voids
all direct search and improves the $n''=n$ record by $\sim47$ orders of magnitude; given a parity
dichotomy, a proof that longer cycles are strictly harder, an Erd\H{o}s--Wintner density theorem, and
a function--field theorem localising the obstruction to the archimedean place; developed an exact
Th\=abit/Euler--type constructive calculus and a bilinear breeder form whose central obstruction we
identify precisely; and extracted the maximal
honest consequence of the construction --- a conditional reduction $(\mathrm{H})\Rightarrow$
\textsc{yes} to an explicit prime--density hypothesis, together with a quantitative ledger showing
that $(\mathrm{H})$ is \emph{not} a feasible finite search (the supply--demand conflict through
$K=GN^2+c^2$ keeps the expected certificate count per reachable frame below $1$). The
rigidity result and the barrier are fully Lean--verified (extremality proved, not assumed), and the
finite verification of Proposition~\ref{prop:close59} sharpens the barrier to $|P\cup Q|\ge60$,
$\min(\prod P,\prod Q)>3.50\times10^{57}$. The heuristics weakly favour \textsc{yes} (existence), hence weakly against
the period--one Ufnarovski--\AA hlander conjecture, but we claim no proof in either direction.

Open problems: (i) formalise the finite verification of Proposition~\ref{prop:close59} (and decide
whether any level beyond $60$ can be closed, e.g.\ by resolving the Pell--type tail family of
Remark~\ref{rem:tier60}); (ii) decide
the parity dichotomy's all--odd case (we expect it empty); (iii) certify more digits of $\delta$ with
an effective error term; (iv) settle the density hypothesis $(\mathrm{H})$ of
Section~\ref{sec:breeder}, or exhibit a non--circular frame family escaping the $\tau(K)/G$ ceiling;
(v) find a $\mathbb{Z}$--side structure detecting the archimedean obstruction of
Section~\ref{sec:ff} --- one that must be non--local, since by Proposition~\ref{prop:nolocal} no single
modulus obstructs; (vi) settle \#307.

\begin{remark}[The place of the problem: a local--global dictionary]\label{rem:place}
The pieces above assemble into a dictionary. \textup{(i)} For a fixed support $U$, the subset form of
Proposition~\ref{prop:subset}, $A(\Sigma-A)=D^2$, satisfies the \emph{Hasse principle as an equation}:
an integer solution $A$ exists iff one exists over $\mathbb{R}$ and over every $\mathbb{Z}_\ell$ (the
discriminant $\Sigma^2-4D^2$ must be a nonnegative square, and squares obey local--global). Its
real--place condition is $\Sigma\ge2D$, i.e.\ $\sigma(U)\ge2$: \emph{the barrier is precisely the
real--place solvability condition of the subset form}. \textup{(ii)} The barrier never needed
primality: pairwise coprime integers have distinct smallest prime factors, so
$\sum 1/a_i\le\sum1/p_i$, and \emph{any} pairwise--coprime support of mass $\ge2$ has at least $59$
elements and product at least $\Pi_{59}$. Coprimality and the ordering of $\mathbb{R}$ carry the wall;
primality enters only through rigidity. \textup{(iii)} What remains unsolved is the \emph{realization}
of the root $A^*=\bigl(\Sigma+\sqrt{\Sigma^2-4D^2}\bigr)/2$ as a sub--multiset sum of the cofactors
$D/p$. Its local conditions are exactly the local anatomy of Lemma~\ref{lem:anatomy}: modulo $\ell^2$
for $\ell\in U$ the subset sums occupy precisely the two cosets
$\ell\mathbb{Z}\cup(D/\ell+\ell\mathbb{Z})$ --- $2\ell$ residues, observed exactly ($6$ of $9$, $10$ of
$25$, $14$ of $49$, $22$ of $121$) --- and these are the same constraints the equation itself forces:
no completion separates the equation from its realization. In the language of integral points, then:
\#307 is everywhere locally soluble, with the barrier as its archimedean height condition, and its
resolution is a \emph{strong approximation} question for the thin, prime--structured set of subset sums
--- a failure mode invisible to Brauer--Manin--type obstructions, consistent with every no--go above
and with the existence of cycles over the neighbouring rings. We record this as a dictionary, not a
method.
\end{remark}

\paragraph{Attribution.}
We are careful to separate prior, concurrent, and present contributions (details in
Section~\ref{sec:prior}). \emph{Prior/concurrent:} the two--cycle reformulation is due to Seva
\cite{Seva323194} (2019) and, independently, Bado \cite{Bado2026}; the bound $|P\cup Q|\ge59$ to
J.~Rosen and the disjointness valuation $v_p(\sum1/p_i)=-1$ to R.~Israel \cite{MO320838}; the source
identification to ``Woett'' \cite{MO320838}. Bado \cite{Bado2026} independently obtained the rigidity
cross--matching, the local anatomy, the parity restriction, a one--sided exact test, approximate
solvability, and the union criterion of Proposition~\ref{prop:pyth}. \emph{Present note:} the
identification of the map with the arithmetic derivative and the two--way bridge to the
Ufnarovski--\AA hlander--Kovi\v{c} literature; the product barrier $2.09\times10^{56}$; the
function--field theorem and archimedean localisation; the constructive deficit analysis and seed
obstruction; the Erd\H{o}s--Wintner density theorem; and the $k$--cycle barriers. Both external
sources \cite{Seva323194,Bado2026} have been examined directly and the overlaps above reflect their
actual contents.

\paragraph{Acknowledgments.}
This work was carried out by the author with substantial assistance from Anthropic's Claude. The
AI assistant contributed to the derivations, the exact and heuristic computations, the Lean~4
formalisation, and the drafting of this note; it was also used to stress-test and attempt to
refute each claim. All statements, proofs, attributions, and the final text were reviewed and are
the responsibility of the author. The author thanks the contributors to the MathOverflow thread
\cite{MO320838} --- in particular R.~Israel and J.~Rosen --- whose observations are credited above.

\section{Note added: prior and concurrent work}\label{sec:prior}

After this note was drafted, two items came to light; having now examined both directly, we record
the overlaps precisely.

\paragraph{Seva (MathOverflow 323194, 2019) \cite{Seva323194}.}
Defines $s(Q)=Q\sum_{p\mid Q}1/p$, derives the valuation property $v_p(s(Q))=v_p(Q)-1$, asks for
$s(P)=Q$, $s(Q)=P$, and states this is Problem~\#307 ``in disguise.'' Priority for the two--cycle
reformulation belongs here. He further asks whether every $s$--trajectory stabilises at $0$ --- the
$s$--analogue of the Ufnarovski--\AA hlander conjecture, posed independently of it. (Note $s\ne n'$
globally, e.g.\ $s(p^\alpha)=p^{\alpha-1}$, but $s=n'$ on squarefree integers, where both two--cycle
questions live.)

\paragraph{Bado (preprint, Aix--Marseille, 2026) \cite{Bado2026}.}
Working with $A(S)=\sum_{p\in S}M(S)/p$ and $M(S)=\prod_{p\in S}p$ (our $N_S,D_S$), Bado independently
obtains: the lowest--terms lemma (his Lemma~3.1, our Lemma~\ref{thm:rigidity}); the forcing identities
$A(P)=M(Q)$, $A(Q)=M(P)$ with $P\cap Q=\varnothing$ (his Thm.~A, our Theorem~\ref{thm:structure}); the
two--cycle reformulation $T(T(P))=P$ (his Thm.~B); the valuation signature (his Thm.~6.1, = Israel's);
the local zero--sum congruences (his Prop.~8.1, our Lemma~\ref{lem:anatomy}); the parity restrictions
(his Prop.~8.2, the parity part of Theorem~\ref{thm:parity}); an exact one--sided test (his Thm.~9.1);
the approximation theorem (his Thm.~10.2, sharper than our greedy version in that all primes exceed a
given bound); and the union discriminant / Pythagorean structure (his Thms.~7.1, 7.4). He cites only
Barbeau, Bloom, and Erd\H{o}s--Graham: crucially, he does \emph{not} identify the map with the
arithmetic derivative, and his note contains none of the present note's $n''=n$ bridge, explicit
barrier constant, function--field theorem, density theorem, $k$--cycle bounds, or constructive
deficit analysis. We regard the overlap on the shared structural layer as strong independent
validation, and import his union criterion with credit:

\begin{proposition}[Single--integer test; derivative form of Bado's Thms.~7.1, 7.4]\label{prop:pyth}
If $(a,b)$ is a two--cycle of the arithmetic derivative, then $N=ab$ is squarefree and satisfies
$N'=a^2+b^2$, whence $N'^2-4N^2=(a^2-b^2)^2$ is a perfect square. Conversely, a squarefree $N$ with
$N'^2-4N^2=\square$ determines a unique unordered factorisation $N=ab$ with $N'=a^2+b^2$, on which the
cycle condition reduces to a single equation (Bado's deviation parameter $t=0$, his Thm.~7.6). This
gives a one--integer necessary test;
computationally, no squarefree $N<2\times10^5$ with $\omega(N)\ge2$ passes it ($0$ of $103{,}596$ ---
consistent with Theorem~\ref{thm:barrier}, which forces $N>4\times10^{112}$).
\end{proposition}
\begin{proof}
$N=ab$ with $a,b$ coprime squarefree gives $N$ squarefree and, by Leibniz with $a'=b,\,b'=a$,
$N'=a'b+ab'=b^2+a^2$; then $N'^2-4N^2=(a^2+b^2)^2-4(ab)^2=(a^2-b^2)^2$. The converse is Bado's
\cite{Bado2026}; the emptiness of the test below $10^{112}$ is upgraded from a sieve to a theorem in
Corollary~\ref{cor:emptytest}.
\end{proof}

The Pythagorean equation $N'=a^2+b^2$ deserves to be isolated from the (harder) cycle condition.

\begin{proposition}[Pair form; the one--equation relaxation]\label{prop:pairform}
For coprime squarefree $a,b$, the single equation $a^2+b^2=(ab)'$ holds if and only if there is an
integer $k$ with $a'=b+ak$ and $b'=a-bk$; a two--cycle is exactly the case $k=0$. We call a solution
of $a^2+b^2=(ab)'$ (any $k$) a \emph{Pythagorean pair}.
\end{proposition}
\begin{proof}
The Leibniz identity $a^2+b^2-(ab)'=a(a-b')-b(a'-b)$ (a polynomial identity, verified symbolically
and on $2000$ random pairs) makes the equation equivalent to $a(a-b')=b(a'-b)$; since $\gcd(a,b)=1$
this forces $a\mid a'-b$ and $b\mid a-b'$ with equal quotients $k$, i.e.\ $a'=b+ak$, $b'=a-bk$.
Setting $k=0$ recovers $a'=b$, $b'=a$. (Equivalently, in $\mathbb{Z}[i]$: \#307 asks for $z=a+bi$ with
coprime coordinates and $ab$ squarefree satisfying $|z|^2=(\operatorname{Im}z^2/2)'$ and $k=0$.)
\end{proof}

\begin{corollary}[The single--integer test is provably empty below $10^{112}$]\label{cor:emptytest}
Every Pythagorean pair satisfies $\sum_{p\mid a}1/p+\sum_{q\mid b}1/q=a/b+b/a\ge2$ (divide
$a^2+b^2=(ab)'$ by $ab$ and apply AM--GM), so its support carries at least $59$ primes and
$N=ab\ge\prod(\text{first }59\text{ primes})>7.9\times10^{112}$. Hence no squarefree $N<10^{112}$ with
$\omega(N)\ge2$ satisfies $N'^2-4N^2=\square$: the empirical ``$0$ of $103{,}596$'' of
Proposition~\ref{prop:pyth} is a theorem, and the barrier of Theorem~\ref{thm:barrier} extends from
two--cycles to the strictly weaker one--equation (Pythagorean) problem.
\end{corollary}

\begin{remark}
(a) As $N'=a^2+b^2$ is a sum of two coprime squares, every odd prime dividing $N'$ is
$\equiv1\pmod4$ --- a cheap union--side pre--filter complementing Bado's congruence filters. (When
exactly one of $a,b$ is even this is the hypotenuse of the primitive triple $(|a^2-b^2|,2ab,a^2+b^2)$.)
(b) Since Seva's $s$ agrees with $n'$ on squarefree integers, and his valuation property forces
two--cycle members squarefree, his question is \emph{exactly} \#307, not merely an analogue.
(c) Speculatively, and recorded only for completeness (no mechanism is known): primitive Pythagorean
triples form the Barning--Hall ternary tree, the orbit structure of a reflection group in
$O(2,1;\mathbb{Z})$; a solution of \#307 would occupy a single node whose generator pair
$(M(P),M(Q))$ carries the full prime--set rigidity.
\end{remark}

The deviation parameter of Proposition~\ref{prop:pairform} organises the one--equation problem into a
short graded family.

\begin{proposition}[The deviation ladder]\label{prop:ladder}
Let $(a,b)$ be a Pythagorean pair with $a<b$ and parameter $k$. Writing
$\sigma(m)=\sum_{p\mid m}1/p$, one has $k=\sigma(a)-b/a=a/b-\sigma(b)\in\mathbb{Z}$, and:
\begin{enumerate}[label=\textup{(\arabic*)},leftmargin=2.2em]
\item $k\le0$, with $k=0$ exactly the \textup{(}open\textup{)} two--cycle case. Indeed $b'=a-bk>0$ and
$b>a\ge1$ force $k\le0$ \textup{(}if $k\ge1$ then $b'\le a-b<0$\textup{)}. Moreover
$|k|<\min(b/a,\sigma(b))\le\sigma(b)\le\log\log b+O(1)$, so the ladder has finitely many rungs --- and,
since $\sigma(b)$ barely exceeds $2$ at the relevant scales, in practice only one or two.
\item \emph{(Parity.)} For odd squarefree $m$, $m'\equiv\omega(m)\pmod2$; for even squarefree $m$, $m'$
is odd. Hence if exactly one member is even, $k\equiv\omega(\text{odd member})\pmod2$; if both are odd,
$\omega(a)\equiv\omega(b)\pmod2$ and $k\equiv\omega(a)+1\pmod2$. The case $k=0$ recovers the parity
dichotomy of Theorem~\ref{thm:parity}.
\item Each rung is the system $a'=b+ak$, $b'=a-bk$, equivalently $(\sigma(a)-k)(\sigma(b)+k)=1$ --- a
shifted copy of the cross--match $st=1$ --- so every rung is exactly as rigid as $k=0$. Since
$\sigma(a)+\sigma(b)=a/b+b/a$ is independent of $k$, Corollary~\ref{cor:emptytest} holds on \emph{every}
rung: the whole ladder is empty below $10^{112}$. (We make no quantitative claim that $k=0$ is easier
than the rest.)
\item On the rung $k=-1$ the equations read $a'=b-a$, $b'=a+b$ and $a'^2+b'^2=2(a^2+b^2)$: the
derivative acts on $z=a+bi$ as the orientation--reversing $\sqrt2$--similarity $z\mapsto(-1+i)\,\bar z$.
\item By uniqueness of the splitting \textup{(}Bado, Thm.~7.5\textup{)}, each squarefree $N$ with
$N'^2-4N^2=\square$ carries, under the convention $a<b$, a well--defined invariant
$k(N)\in\mathbb{Z}_{\le0}$; \#307 holds iff $k(N)=0$ for some $N$, and by~(3) the invariant lives on a
domain that is empty below $10^{112}$.
\end{enumerate}
\end{proposition}

The Pythagorean layer behaves very differently over function fields, and refines into two arithmetic
sub--layers over $\mathbb{Z}$.

\begin{proposition}[The Pythagorean layer over function fields]\label{prop:ff-pyth}
Over $\mathbb{F}_q[t]$ with $q$ odd, $a^2+b^2=(ab)'$ has no coprime squarefree monic solutions:
$\deg(a^2+b^2)=2\max(\deg a,\deg b)>\deg a+\deg b-1\ge\deg((ab)')$, the leading coefficient $1$ or
$2\ne0$ preventing cancellation. Over $\mathbb{F}_2[t]$ the layer is nonempty: up to degree $8$
(exhaustively, $23{,}548$ coprime squarefree pairs) its only solutions are the three pairs of the
Artin--Schreier tower $x_{j+1}=x_j(x_j+1)$, namely $(t,t+1)$, $(t^2+t,t^2+t+1)$, $(t^4+t,t^4+t+1)$;
each has $a'=(a+1)'=1$ and $(a(a+1))'=(a+1)+a=1$, and the climb halts at level $4$ (there
$(x_3+1)'=t^2+t\ne1$). These pairs lie off the cycle stratum $k=0$, on the side the integer
normalisation of Proposition~\ref{prop:ladder}(1) forbids. Cycles ($k=0$) remain impossible over every
$\mathbb{F}_q[t]$ (Theorem~\ref{thm:ff}): the function--field world separates the Pythagorean layer
from the cycle layer, sharpening the integer obstruction to ``the archimedean order forces $k\le0$.''
\end{proposition}

\begin{proposition}[Deviation calculus for $k$--cycles]\label{prop:kdev}
For pairwise coprime squarefree $a_1,\dots,a_k$ with product $N$ and deviations
$\kappa_i=(a_i'-a_{i+1})/a_i$ (indices mod $k$): the ratios $a_{i+1}/a_i$ have product $1$
(telescoping), and $\sigma(N)=\sum_i a_{i+1}/a_i+\sum_i\kappa_i$. On the relaxed layer
$\sum_i\kappa_i=0$ --- which contains the cycles $\kappa_1=\dots=\kappa_k=0$ --- AM--GM gives
$\sigma(N)\ge k$, so the distinct--prime barrier $m_k$ of Proposition~\ref{prop:kcycles} extends to the
relaxed problem for every $k$, and to genuine $k$--cycles for $k\le3$; for $k\ge4$ it inherits the
pairwise--coprimality caveat of Proposition~\ref{prop:kcycles} (the bound then holds for the multiset
reciprocal sum). At $k=2$, $\kappa_1=-\kappa_2$ is the deviation of Proposition~\ref{prop:ladder} and
$\sum\kappa_i=0$ is $a^2+b^2=(ab)'$.
\end{proposition}

\begin{proposition}[Split discriminant]\label{prop:split}
For coprime squarefree $a,b$ with $N=ab$,
\[ N'+2N-(a+b)^2=N'-2N-(a-b)^2=(ab)'-a^2-b^2. \]
Hence $N$ is the product of a Pythagorean pair iff $N'-2N$ and $N'+2N$ are \emph{both} perfect squares
$d^2,s^2$ (their product is Bado's discriminant $N'^2-4N^2$), with $a,b=(s\pm d)/2$. The minus factor
carries the obstruction: $N'-2N=(a-b)^2\ge0$ forces $\sigma(N)\ge2$, hence the $59$--prime barrier ---
indeed $0$ of $103{,}596$ squarefree $N<2\times10^5$ satisfy $N'-2N=\square$
(Corollary~\ref{cor:emptytest} in one line). The plus factor is satisfiable at small scale ($114$ such
$N<2\times10^5$). In characteristic $2$ the obstructing factor vanishes ($a-b=a+b$) and the criterion
degenerates to $(a+b)^2=N'$: the permissiveness of Proposition~\ref{prop:ff-pyth} and the barrier of
Corollary~\ref{cor:emptytest} are the absence and presence of one and the same algebraic factor.
\end{proposition}

\begin{proposition}[Structure of the plus--layer]\label{prop:plus}
Call squarefree $N$ with $\omega(N)\ge2$ a \emph{plus--hit} if $N'+2N$ is a perfect square; the first
are $130,145,305,689,1745,1870,1970,2353,\dots$ \textup{(}$114$ below $2\times10^5$\textup{)}.
\textup{(i)} For a semiprime $N=pq$, $2(N'+2N)+1=(2p+1)(2q+1)$, so a plus--hit forces
$(2s)^2\equiv-2\pmod\ell$ for every prime $\ell\mid(2p+1)(2q+1)$, whence $\ell\equiv1,3\pmod8$ and
$p\equiv q\equiv1\pmod4$ (verified on all $59$ semiprime instances; this explains the dominance of $5$
and $13$). \textup{(ii)} \textup{[}empirical\textup{]} $\#\{\text{plus--hits}\le X\}\approx0.25\sqrt X$
(ratios $0.250,0.253,0.255$ at $X=5\!\times\!10^4,10^5,2\!\times\!10^5$): the plus factor behaves like
a random integer of its size, so --- with Proposition~\ref{prop:split} --- the whole integer
obstruction sits in the minus factor. \textup{(iii)} Fixing $p=5$, $q=(s^2-5)/11$ with $s$ even and
$s\equiv\pm4\pmod{11}$ gives $q\equiv1\pmod4$, and each prime value yields the plus--hit $N=5q$
($68$ with $s<2000$, up to $N=1{,}774{,}805$); Bunyakovsky's conjecture for this quadratic thus
implies infinitely many plus--hits. The two discriminant factors separate cleanly: the plus--layer is
Bunyakovsky--class, the minus--layer is the barrier.
\textup{(iv)} [Remark] The full cycle condition states that $(d^2,N',s^2)$ is a three--term
arithmetic progression with common difference $2N$ whose outer terms are squares --- a
congruent--number--shaped configuration; the elliptic--curve machinery of that theory does not
transfer because $N'$ is not an algebraic function of $N$. A $251$--term b--file to $N<10^6$
(counting ratio $0.251$) and an OEIS submission draft accompany this note; the plus--hit
sequence matched no OEIS entry in our search.
\textup{(v)} [P, C] \emph{General local law (QR tournament).} For every odd prime $\ell\mid N$ of
a plus--hit, $N'+2N\equiv N/\ell\pmod\ell$ (all cofactor terms $N/p$ with $p\ne\ell$ carry a factor
of $\ell$; only $N/\ell$ survives), so $N/\ell$ is a nonzero quadratic residue mod $\ell$. Verified
with zero violations on all $114$ plus--hits below $2\times10^5$ and all $251$ terms to $10^6$;
the semiprime $p\equiv q\equiv1\pmod4$ theorem of~\textup{(i)} is its $\omega=2$ shadow via quadratic
reciprocity.
\end{proposition}

\begin{corollary}[Quadratic--residue filter for \#307 solutions]\label{cor:qr}
Let $(P,Q)$ solve Erd\H{o}s~\#307 and set $a=D_P$, $b=D_Q$, $N=ab$. Since $N'=a^2+b^2$ and
hence $N'+2N=(a+b)^2$ automatically, every prime $\ell\in P\cup Q$ satisfies $(N/\ell\mid\ell)=+1$:
the cofactor $N/\ell$ is a nonzero quadratic residue modulo $\ell$. Indeed $\ell$ divides exactly
one of $a,b$ (by $P\cap Q=\varnothing$), so $\ell\nmid a+b$ and $(a+b)^2\not\equiv0\pmod\ell$;
combined with the identity $N'+2N\equiv N/\ell\pmod\ell$ of~\textup{(v)}, this forces $N/\ell$ to
be a nonzero square mod $\ell$. As a search filter the condition prunes candidate unions by a factor
${\le}2^{-|P\cup Q|}\le2^{-59}$ (heuristically each of the $\ge59$ primes imposes an independent
QR constraint); for genuine cycles it is automatic and does not alter the existence heuristics.
We note that Corollary~\ref{cor:qr} is \emph{not} independent of Bado: it is the squared form of his
congruence $D_Q\equiv D_P/\ell\pmod\ell$ (itself $N_P$ reduced mod $\ell$ by cofactor survival),
giving $N/\ell\equiv D_Q^2\pmod\ell$. Squaring is lossy --- $r$ and $-r$ share a square --- so Bado's
linear congruence is strictly stronger, pinning $D_Q\bmod\ell$ exactly where the residue filter
determines it only up to sign; the filter's one added value is being \emph{split--free} (the residue
$N/\ell\bmod\ell$ is computable from the candidate union alone, with no knowledge of the $P\mid Q$
partition). The general law of Proposition~\ref{prop:plus}(v), by contrast, holds for \emph{any}
squarefree plus--hit with no $P\mid Q$ split and is not subsumed by Bado.
\end{corollary}

The plus-- and minus--hits both organise into one--prime ``slot'' families, which makes the structure
of Proposition~\ref{prop:split} explicit at every $\omega$.

\begin{proposition}[Slot calculus for the discriminant layers]\label{prop:slot}
Let $N_0$ be squarefree, $r\nmid N_0$ a prime, $N=N_0r$, and $E_0\coloneqq N_0'+2N_0$. By Leibniz
($N'=N_0'r+N_0$),
\[ N'+2N=rE_0+N_0,\qquad N'-2N=r(N_0'-2N_0)+N_0, \]
so $N$ is a plus--hit iff $r=(s^2-N_0)/E_0$ for some integer $s$ (equivalently $s^2\equiv N_0\pmod{E_0}$).
\begin{enumerate}[label=\textup{(\alph*)},leftmargin=2.2em]
\item \emph{Every $\omega$--stratum.} A plus--hit $N$ ($\omega\ge2$) factors, through any prime divisor
$r$, as $N=N_0r$ with $N_0=N/r$ squarefree, and lies in the quadratic family
$\{N_0(s^2-N_0)/E_0\}_s$ of its base --- generalising the $p=5$ family of
Proposition~\ref{prop:plus}(iii) beyond semiprimes (each step adds one prime,
$\omega(N)=\omega(N_0)+1$). The congruence $s^2\equiv N_0\pmod{E_0}$ is solvable for $669$ of the
$3041$ squarefree $N_0\in[2,5000]$ ($22.0\%$); solvability is necessary but not sufficient for the
family to contain a hit.
\item \emph{Recursive towers: a trichotomy.} When the base is itself a plus--hit, $E_0=s_0^2$ is a
square; below $2\times10^5$ the $114$ plus--hit bases split into three kinds. \emph{No residue} ($84$):
$s^2\equiv N_0\pmod{E_0}$ is unsolvable and the family is empty --- a quadratic--residue obstruction,
occurring for both even $E_0$ ($N_0=305$, with $305$ a non--residue mod $13$) and odd $E_0$
($N_0=1870$, $E_0=73^2$). \emph{Fixed divisor} ($13$, each with fixed divisor $2$): the family is
nonempty but every slot is even, so only $r=2$ could serve and never gives a square --- sterile (e.g.\
$N_0=145$, $E_0=18^2$). \emph{Live} ($17$): fixed divisor $1$ and a prime slot is exhibited (e.g.\
$N_0=130$, first prime slot $r=83399$, giving the plus--hit $N=10{,}841{,}870$ with $N'+2N=5487^2$).
\item \emph{Minus layer; the barrier in slot form.} For $\sigma(N_0)<2$ the minus identity gives
$N'-2N\ge0$ iff $r\le1/(2-\sigma(N_0))$, i.e.\ iff $\sigma(N)\ge2$: this is the slot--variable form of
Corollary~\ref{cor:emptytest}, not an independent proof of it. The bound $1/(2-\sigma(N_0))$ is
unbounded as $\sigma(N_0)\to2^-$ --- over squarefree $N_0<10^5$ its maximum is only $1.524$ (at
$N_0=30030$), so no such base admits a prime minus--slot, yet for the $58$--prime primorial the bound
exceeds $793$ and the next prime $277$ fits, producing the $59$--prime primorial, the smallest
minus--hit base. A prime minus--slot thus first becomes feasible \emph{exactly} at the $59$--prime
barrier; the minus--empty--below--barrier statement rests on the primorial growth of
Corollary~\ref{cor:emptytest}, not on the figure $1.524$.
\end{enumerate}
\end{proposition}
\noindent (The two identities are checked symbolically and on $2\times10^4$ random pairs; the base
trichotomy, slot lists, and the $1.524$ maximum by exhaustive enumeration.)

The two square conditions over a base combine into a single binary--form representation, which places
the obstruction inside Gauss's genus theory.

\begin{proposition}[The Pythagorean condition as a binary--form representation]\label{prop:form}
For squarefree $N_0$ and prime $r\nmid N_0$, $N=N_0r$, the identity
\[ (2N_0-N_0')(N'+2N)+(2N_0+N_0')(N'-2N)=4N_0^2 \]
holds unconditionally (Leibniz; checked on $2\times10^4$ random pairs). Hence if $N$ is a Pythagorean
pair product over $N_0$ --- both $N'+2N=s^2$ and $N'-2N=d^2$ squares --- then $(s,d)$ represents
$4N_0^2$ by the binary form $Q_{N_0}(s,d)=A\,s^2+B\,d^2$ with $A=2N_0-N_0'=N_0(2-\sigma(N_0))$,
$B=2N_0+N_0'=N_0(2+\sigma(N_0))$, of discriminant $-4AB=4(N_0'^2-4N_0^2)$; conversely such a
representation yields a Pythagorean pair over $N_0$ exactly when the recovered slot $r=(s^2-N_0)/E_0$
is a positive integer, prime, and coprime to $N_0$ (the integrality $s^2\equiv N_0\pmod{E_0}$ holds on
every representation we found, $N_0<2\times10^4$). For $\sigma(N_0)<2$ the form is positive--definite
($A>0$), and below the barrier representations are rare and all ``ghosts'': the only squarefree
$N_0\le3000$ representing $4N_0^2$ are $N_0=30$ (with $(s,d)=(11,1)$) and $N_0=1722$ ($(83,1)$), each
carrying the non--admissible slot $r=1$, and no live prime slot occurs. As $\sigma(N_0)\to2^-$ the
leading coefficient $A=N_0(2-\sigma(N_0))\to0$ and the ellipse $A\,s^2+B\,d^2=4N_0^2$ degenerates to an
arbitrarily thin sliver: the geometric face of the mass--2 barrier. Consequently \#307 is the existence
of a base with $\sigma(N_0)\ge2-1/r$ whose form $Q_{N_0}$ represents $4N_0^2$ with an admissible,
integral, \emph{prime} slot and deviation $k=0$.
\end{proposition}
\noindent (We make no class--number claim: each base yields a form of a \emph{different} discriminant,
and the question is single--value representation, not a class count; only the degeneration $A\to0$ is
used.)

\begin{proposition}[A principal--genus necessary condition, and a cascade on solutions]\label{prop:genus}
As $\gcd(N_0,N_0')=1$, the value $4N_0^2$ is a square coprime to the odd part of $\operatorname{disc}
Q_{N_0}$, so every genus character is $+1$ on it and a primitive form represents it only from the
principal genus \cite{Cox}. Reading the odd assigned characters off $Q_{N_0}(1,1)=4N_0$ gives the
computable necessary condition
\[ (N_0\mid\ell)=+1 \quad\text{for every odd prime }\ell\mid (N_0'^2-4N_0^2)\text{ to odd multiplicity} \]
(a $2$--adic condition accompanies it; for odd $N_0$ of even $\omega$ the form is imprimitive and must be
primitivised first --- so we state the criterion on the odd part, with a $2$--adic caveat). Both ghost
bases $30,1722$ satisfy it, and about $16\%$ of squarefree bases pass in the ranges tested (an
\emph{upper bound}, using only the odd characters, and mildly range--dependent). \emph{Cascade.} In a
solution of \#307 with $N=D_PD_Q$, each $\ell\in P\cup Q$ gives a cofactor $N/\ell$ realising the
Pythagorean structure of $N$ over base $N/\ell$; every cofactor with $\sigma(N/\ell)<2$ --- all of them
for a minimal first--$59$--primes solution --- must then satisfy the principal--genus criterion, a
cascade of $\ge59$ computable genus conditions complementing the quadratic--residue filter of
Corollary~\ref{cor:qr}. \textup{(The further class--group and slot--admissibility stages of the
fertility gap we observe only numerically at toy scale.)}
\end{proposition}

Below the barrier the form has exactly one source of solutions, and it is entirely classical.

\begin{proposition}[The sub--barrier ghosts are pronic Giuga numbers]\label{prop:giuga}
Every representation of $4N_0^2$ by $Q_{N_0}$ over squarefree $N_0\le3\times10^4$ is a ghost
$(s,d)=(\sqrt{4N_0+1},1)$, occurring exactly when $N_0'=N_0+1$ with $N_0$ pronic; in range only
$30=5\cdot6$ and $1722=41\cdot42$. The two lines through such bases are classical: since $n'=n\,\sigma(n)$
for squarefree $n$,
\[ \{n:n'=n-1\}=\text{primary pseudoperfect numbers }(\textstyle\sum_{p\mid n}\tfrac1p+\tfrac1n=1),\qquad
   \{n:n'=n+1\}=\text{Giuga numbers (quotient one)}. \]
The derivative characterisations are due to Lava (OEIS \cite{OEIS_A054377,OEIS_A007850}, 2009 --- the
Giuga form $n'=n+1$ being Lava's \emph{conjecture}, with the generalisation $n'=an+1$, $a\in\NN$, the
theorem of Grau--Oller-Marc\'en \cite{GrauOllerMarcen2012}); the underlying number classes are
Butske--Jaje--Mayernik \cite{BJM2000} (exactly one primary pseudoperfect number per $\omega\le8$, eight
known) and Giuga \cite{Giuga1950,BBBG1996}; and the oblong chain ($N$ with $N+1$ prime breeds the next
term $N(N+1)$) is classical in both directions (Forgues and Sondow, OEIS \cite{OEIS_A054377}). To
$2\times10^6$ the lines are $\{n'=n-1\}=\{2,6,42,1806,47058\}$ and $\{n'=n+1\}=\{30,858,1722,66198\}$,
and the chain correctly predicts the next terms: $47059$ is prime, so $47058\cdot47059=2{,}214{,}502{,}422$
is the sixth primary pseudoperfect number, and $47057$ is prime, so the pronic
$47057\cdot47058=2{,}214{,}408{,}306$ is the fifth Giuga number --- both already known, recovered here by
structure. Our contribution is the reframing: the chain step is the one--line Leibniz identity
$N'=N-1\Rightarrow(N(N+1))'=N'(N+1)+N^2=N(N+1)-1$, exhibiting the chain as orbits of the Artin--Schreier
map $x\mapsto x(x+1)=x^2+x$ --- the same map, with the same halt ($2\to6\to42\to1806$, stopping at
$1807=13\cdot139$), as the $\mathbb{F}_2[t]$ tower of Proposition~\ref{prop:ff-pyth}: characteristic $2$
and $\mathbb{Z}$ run identical breeders on the two lines.
\end{proposition}

\begin{remark}[Converse, and the bridge]
Every pronic solution of $n'=n+1$ is $m(m+1)$ with $m$ prime and $m+1$ primary pseudoperfect: Leibniz
gives $m'(m+1)+m(m+1)'=m^2+m+1$, and writing $m'=km+1$ forces $(m+1)'=(1-k)m-k$, which is $\ge0$ only for
$k=0$, whence $m'=1$ ($m$ prime) and $(m+1)'=m$. (The congruence $m'\equiv1\pmod m$ alone does \emph{not}
force $m$ prime --- its composite solutions are precisely the Giuga numbers $30,858,1722,66198,\dots$ ---
so primality comes from the sign elimination $k=0$.) Thus the only sub--barrier life in the form equation
is the classical Giuga/primary--pseudoperfect world, known to sit in a web with Sylvester's sequence,
Zn\'am's problem and the Erd\H{o}s--Moser equation \cite{SondowMacMillan2017} and with the arithmetic
derivative \cite{GrauOllerMarcen2012,GrauOllerMarcenSadornil2021}; we claim none of this web, only that
\#307 meets it through $n'=n\,\sigma(n)$, and that its $a\ge2$ (equivalently $\sigma>2$) branch lands on
the same $\ge59$--prime barrier \cite{BorweinWong1997} that bars genuine cycles
(Corollary~\ref{cor:emptytest}).
\end{remark}

Writing both members of a cycle as members of these lines stratifies \#307 by the gap.

\begin{proposition}[Gap stratification: \#307 as adjacency of $\mu$--Sondow lines]\label{prop:gap}
For a squarefree two--cycle $a'=b$, $b'=a$ with $a>b$ and gap $d=a-b$, the smaller member lies on the line
$n'=n+d$ and the larger on $n'=n-d$ (as $b'-b=d$ and $a'-a=-d$), with $\gcd(d,ab)=1$ (a prime dividing $d$
and $a$ would divide $b=a-d$); conversely $x'=x+d$ with $(x+d)'=x$ is a two--cycle. These opposite--sign
lines are the quotient--one $\mu$--Sondow numbers of parameter $\mp d$ \cite{GrauOllerMarcenSadornil2021}.
Hence \#307 holds iff for some $d\ge1$ the lines $n'=n+d$ and $n'=n-d$ carry members exactly $d$ apart; for
$d=1$ this is \emph{a Giuga number adjacent to a primary pseudoperfect number}. By the parity dichotomy
(Theorem~\ref{thm:parity}) the mixed--parity case forces $d$ odd and the both--odd case ($d$ even) carries
the $10^{2519}$ bound, so every cycle below that bound has odd gap and $d=1$ is the lead stratum. A $d=1$
pair $(k,k+1)$ has exactly one odd member, hence requires either an odd Giuga number (known to have at least
$9$ prime factors \cite{BBBG1996,BorweinWong1997}) or an odd primary pseudoperfect number (none is known):
the lead stratum inherits two classical open problems. Tabulating cycle--eligible members (squarefree,
coprime to $d$) to $2\times10^6$ for all $d\le50$ gives $21$ on the plus lines (none of even gap, consistent
with the parity remark) and $81$ on the minus lines, with \emph{no} adjacent pair at any $d\le50$ --- the
barrier in gap--stratified form. \textup{(}For coprime $M$ on $n'=n-\alpha$ and $N$ on $n'=n+\beta$, Leibniz
gives $(MN)'-2MN=\beta M-\alpha N$, so $\beta M-\alpha N=\square$ would manufacture a minus--side square in
the sense of Proposition~\ref{prop:split}; the case $(\alpha,\beta)=(1,1)$, $M-N=1$ is the $d=1$ cycle
itself, and the known finite Giuga/primary--pseudoperfect lists contain no coprime pair, all members being
even.\textup{)} To our knowledge this gap--stratified equivalence and the pairing of opposite--sign lines at
fixed additive distance are new; the lines themselves are \cite{GrauOllerMarcenSadornil2021}.
\end{proposition}

The line parameter $\delta(n)=n'-n$ obeys an affine recursion that organises the classical members into a
genealogy.

\begin{proposition}[The $\delta$--recursion and the genealogy of the classical lines]\label{prop:delta}
With $\delta(n)=n'-n$, for squarefree $M$ and a prime $p\nmid M$,
\[ \delta(Mp)=p\,\delta(M)+M \qquad(\text{Leibniz; checked on }10^3\text{ random pairs}). \]
Solving $p\,\delta(M)+M=\pm d$ gives the child rule $p=(M\mp d)/e$ from a parent $M$ on the minus line
$\delta(M)=-e$. Every known Giuga number arises this way from a small--parameter parent: $30=6\cdot5$
($e=1$), $858=66\cdot13$ ($e=5$), $1722=42\cdot41$ ($e=1$), $66198=1122\cdot59$ ($e=19$),
$2{,}214{,}408{,}306=47058\cdot47057$ ($e=1$); the seed $47058$ itself resolves as $6\to66\to1518\to47058$
(with $\delta=-1,-5,-49,-1$), so the previously unexplained members $858$ and $66198$ acquire genealogies.
Two members of a cycle can never share a parent: siblings $M(M\mp d)/e$ differ by $2Md/e>d$, since
$e=M-M'<2M$. The oblong case $e=1$ is the classical chain of Proposition~\ref{prop:giuga} (its
$\mu$--Sondow form is \cite{GrauOllerMarcenSadornil2021}); the recursion, the general--$e$ genealogy, and the
resulting small--$|\delta|$ tree do not appear there.
\end{proposition}

\begin{corollary}[Nonemptiness families for the $\mu$--Sondow conjecture]\label{cor:gos}
In the notation of \cite{GrauOllerMarcenSadornil2021} (their $S_\mu$; the quotient--one slice is the line
$n'=n-\mu$), the recursion produces explicit members exceeding $|\mu|$ in three infinite parameter families:
$2p\in S_{p-2}$ for every odd prime $p$; $Mp\in S_{p-M}$ for every primary pseudoperfect number $M$ and
prime $p\nmid M$; and $Gp\in S_{-(p+G)}$ for every Giuga number $G$ and prime $p\nmid G$ --- in each case
$\mu/n+\sum_{q\mid n}1/q=1$ exactly (checked symbolically and on $151$ instances). This settles their
nonemptiness conjecture ($S_\mu\cap(|\mu|,\infty)\ne\varnothing$, their Conjecture~1(ii)) for infinitely
many parameters, including infinitely many outside their verified range $-145<\mu<673$; these families do
not appear among their constructions, which scale by $\operatorname{rad}(\mu)$ rather than extend by a new
prime. Their two open instances resist precisely because every quotient--one parent route hits a composite:
for $\mu=-145$, $M-145$ is composite for all eight known primary pseudoperfect numbers
($1661=11\cdot151$, $46913=43\cdot1091$, and the three larger cases, including the recently found eighth
member) while $p+G=145$ has no prime solution; for $\mu=673$, $M+673$ is composite for all eight
($679=7\cdot97$, $2479=37\cdot67$, $47731=59\cdot809$, the rest ending in $5$) and $p=675$ fails. A prime
in either family at the next generation would settle a published open instance; an exhaustive search found
none. Concretely: no quotient--one member of $S_{-145}$ or $S_{673}$ exists with \emph{any} prime--factor
cofactor $\le10^{11}$ (every squarefree parent $M\le10^{11}$ was scanned against the child rules, a region
whose members reach ${\sim}10^{22}$, together with a depth--two descent), the quotient--zero slice of
$S_{-145}$ is empty outright ($n'=145$ forces $n\le5256$ by $n'\ge2\sqrt n$, and the finite check is
empty), and quotients $\ge2$ force $\sigma\gtrsim2$, hence at least $59$ prime factors. Both instances
therefore remain open, with the quotient--one frontier moved from the $10^{10}$ of
\cite{GrauOllerMarcenSadornil2021} to cofactor scale $10^{11}$.
\end{corollary}

\begin{remark}[The look--back]
Along small--$|\delta|$ paths the forced prime is $p\approx M/e$, so members grow doubly exponentially and
the $10^{56}$--$10^{113}$ wall zone lies only a few generations from the seeds: a sparse, recursively
enumerable tree (a child needs $(M-j)/e$ prime for small $j$). The adjacency of two such trees at a fixed
additive distance places the final condition in the Pillai/$S$--unit genre --- two multiplicatively
generated sequences meeting at a fixed additive distance --- where fixed--support finiteness is classical but
growing--support statements are conjectural; a lens, not a transfer. Re--expressed in $\delta$--coordinates
the obstruction is once more $\sigma>1$ (positivity of $\delta$) plus integer reachability $\delta=\pm d$
under forced primes: the same additive--multiplicative exactness reappears in every coordinate tried (mass,
the deviation $k$, the form $Q_{N_0}$, genus theory, $\mu$--Sondow lines, $\delta$--dynamics). The reframings
buy structure, not a reduction.
\end{remark}

\begin{proposition}[Even--gap plus lines are empty below $3.23\times10^9$]\label{prop:oddthr}
An even squarefree $x=2u$ has $\delta(x)=2u'-u$ odd, so every line member of even $\delta$ is odd (with
$\omega$ odd). An odd member of a plus line has $\sigma>1$ over odd primes, and the least odd squarefree
integer of mass exceeding $1$ is $3\cdot5\cdots29=3{,}234{,}846{,}615$ (the $8$--prime product
$3\cdots23=111{,}546{,}435$ has mass $0.99896$). Hence every plus line of even gap $d$ is empty below
$3{,}234{,}846{,}615$, proving the census observation (no even--gap plus member to $2\times10^7$) and placing
an explicit floor under the smaller member of any both--odd configuration.
\end{proposition}

The lines also admit parity, infinitude, and congruence refinements.

\begin{proposition}[Parity floor for odd line members]\label{prop:parityfloor}
For odd squarefree $n$ every cofactor $n/p$ is odd, so $n'\equiv\omega(n)\pmod2$. Hence an odd primary
pseudoperfect number ($n'=n-1$, even) and an odd Giuga number ($n'=n+1$, even) both have $\omega(n)$ even;
combined with the classical bound $\omega\ge9$ \cite{BBBG1996,BorweinWong1997} this forces $\omega\ge10$, a
parity sharpening (for the primary pseudoperfect side, the eight smallest odd primes give $\sum1/p=0.99896<1$,
so $\omega\ge9$ directly). The floor cascades when small primes are avoided: for an odd squarefree member of
any plus line, or of a quotient--one line (where $\sigma(n)\ge1-1/n$), exact greedy/exchange computation
gives $\omega\ge27$ if $3\nmid n$; $\omega\ge66$ if $3,5\nmid n$; $\omega\ge144$ if $3,5,7\nmid n$; and
$\omega\ge16$ if $3\mid n$, $5\nmid n$ --- each avoided small prime multiplies the required complexity, and
for odd Giuga or primary pseudoperfect numbers the parity constraint rounds odd floors up (e.g.\ $3\nmid n$
forces $\omega\ge28$). (The restriction to plus or quotient--one lines is essential: on a general minus line
$\sigma$ can be small --- every prime $p$ trivially satisfies $p'=p-(p-1)$ --- so no such floor exists
there.)
\end{proposition}

\begin{proposition}[Infinitely many gaps carry a nonempty plus line]\label{prop:plusinf}
By the recursion of Proposition~\ref{prop:delta}, $\delta(30p)=p\,\delta(30)+30=p+30$ for every prime $p>5$
(as $\delta(30)=1$), so the line $n'=n+(p+30)$ is nonempty; by Euclid infinitely many distinct gaps $d$ carry
a nonempty plus line, and the same runs from any Giuga number $G$ ($\delta(Gp)=p+G$). This fixes the census
tail: $330,390,510,570=30\cdot\{11,13,17,19\}$ lie on the gaps $d=41,43,47,49$. (Per--gap infinitude remains
open, of Bunyakovsky type.)
\end{proposition}

\begin{proposition}[A difference--$36$ congruence modulo $72$]\label{prop:mod72}
Let $n=6m$ be a Giuga or primary pseudoperfect number with $\gcd(m,6)=1$, $n'=n+\varepsilon$
($\varepsilon=+1$ Giuga, $-1$ primary pseudoperfect). Then $6m'=m+\varepsilon$, and with
$m'\equiv\omega(m)\pmod2$ this yields $n\equiv36\,\omega(n)-78\pmod{72}$ on the Giuga side and
$n\equiv36\,\omega(n)-66\pmod{72}$ on the primary pseudoperfect side; in particular equal $\omega$ forces
equal residue. The congruence is verified on all twelve Giuga numbers and all seven $6$--divisible primary
pseudoperfect numbers listed in the OEIS (we recomputed the factorisations of the five largest Giuga
members; all are squarefree, with $\omega=7,7,8,8,8$), and gives, at the modulus $72$, the
difference--$36$ arithmetic progression observed by Sondow--MacMillan \cite{SondowMacMillan2017}; their
finer modulus--$288$ pattern remains empirical.
\end{proposition}

The modulus--$72$ argument is the $D=6$ case of a general transport of the line equation across a fixed
divisor.

\begin{proposition}[Divisor transport]\label{prop:transport}
Let $n=Dm$ with $D,m$ coprime squarefree and $n'=n+\varepsilon$ ($\varepsilon=\pm1$) a quotient--one member.
Leibniz transports the line equation to
\[ D\,m'=(D-D')\,m+\varepsilon. \]
\begin{enumerate}[label=\textup{(\alph*)},leftmargin=2.2em]
\item If $\sigma(D)>1$ (so $D'>D$; least case $D=30$), the right side is negative for $m>1$ while the left
is nonnegative; the only escape is $m=1$, $\varepsilon=+1$. Hence \emph{no primary pseudoperfect number is
divisible by $30$, and the only Giuga number divisible by $30$ is $30$ itself} (elementary; we have not
located it in the literature) --- consistent with the full known corpus, of which only $30$ carries the
factor~$5$.
\item At a primary pseudoperfect block $D$ (where $D-D'=1$) the slice $m'=1$ reads $m=D-\varepsilon$ prime:
the two classical chain rules --- $D(D+1)$ primary pseudoperfect when $D+1$ is prime, $D(D-1)$ Giuga when
$D-1$ is prime \cite{GrauOllerMarcenSadornil2021,SondowMacMillan2017} --- are the two signs of one
transported equation ($1806=42\cdot43$ and $1722=42\cdot41$ sit on the same slice $D=42$).
\item For any gap $d$ the same transport stratifies the line $n'=n+d$ by divisor blocks $D$ into the steeper
lines $D\,m'=(D-D')m+d$, a further enumeration coordinate for the genealogy of
Proposition~\ref{prop:delta}.
\end{enumerate}
\end{proposition}

Finally, the form obstruction can be averaged: it eliminates almost every integer, though not the thin
stratum where a solution could live.

\begin{proposition}[Average rarity of form--barrier survivors]\label{prop:survivors}
Call a squarefree $N$ \emph{blocked} at an odd prime $\ell$ if $\ell\nmid N$, $N'\equiv\pm2N\pmod\ell$, and
$(N\mid\ell)=-1$, and a \emph{survivor} if it is blocked at no odd prime. A blocked $N$ admits no
representation $E\,s^2+D\,d^2=4N^2$ (reducing mod $\ell$ forces $N$ to be a quadratic residue, against
$(N\mid\ell)=-1$; this is the genus fragment of Proposition~\ref{prop:genus}), so every union of a \#307
solution is a survivor. Now $N'\equiv N\,S_\ell(N)\pmod\ell$ with $S_\ell(N)=\sum_{p\mid N}p^{-1}$ an
\emph{additive} function, so blocking at $\ell$ is an additive--function event ($S_\ell\equiv\pm2$) twisted by
the quadratic character, of limiting density $1/(\ell+1)$. As $\sum_\ell1/(\ell+1)$ diverges (Mertens), the
Selberg--Delange/Halász theory of additive functions in residue classes \cite{Tenenbaum} gives that the
survivors have natural density zero among the squarefree integers (unconditional), with
$\#\{\text{survivors}\le X\}=O\!\bigl(X/\log\log\log X\bigr)$ under a standard --- but here unproved ---
uniformity estimate. Thus the single--prime form obstruction removes \emph{almost every} squarefree integer.
We stress the scope: survivors remain \emph{infinite}, and a \#307 solution --- one integer beyond
$2.09\times10^{56}$, in the density--zero stratum $\sigma(N)\ge2$ --- is untouched by such an average
statement, which speaks only to the typical integer, not to existence.
\end{proposition}

\begin{thebibliography}{9}
\bibitem{Barbeau1961} E.~J.~Barbeau, \emph{Remarks on an arithmetic derivative}, Canad.\ Math.\
Bull.\ \textbf{4} (1961), 117--122.
\bibitem{Barbeau1976} E.~J.~Barbeau, \emph{Computer challenge corner: Problem 477}, J.\ Recreational
Math.\ \textbf{9} (1976/77), 30.
\bibitem{ErdosGraham1980} P.~Erd\H{o}s and R.~L.~Graham, \emph{Old and New Problems and Results in
Combinatorial Number Theory}, Monogr.\ Enseign.\ Math.\ \textbf{28}, Geneva, 1980, p.~38.
\bibitem{GrahamEgyptian} R.~L.~Graham, \emph{Paul Erd\H{o}s and Egyptian fractions}, in \emph{Erd\H{o}s
Centennial}, Bolyai Soc.\ Math.\ Stud.\ \textbf{25}, Springer, 2013, pp.~289--309.
\bibitem{Watanabe2020} T.~Watanabe, \emph{New examples of the representation of $1$ by the sum of
reciprocals of semiprime numbers}, arXiv:2009.03275 (2020).
\bibitem{BloomShifted} T.~F.~Bloom, \emph{Unit fractions with shifted prime denominators}, Proc.\
Roy.\ Soc.\ Edinburgh Sect.\ A (2024); arXiv:2305.02689.
\bibitem{UfnarovskiAhlander2003} V.~Ufnarovski and B.~\AA hlander, \emph{How to differentiate a
number}, J.\ Integer Seq.\ \textbf{6} (2003), Article 03.3.4.
\bibitem{Kovic2012} J.~Kovi\v{c}, \emph{The arithmetic derivative and antiderivative}, J.\ Integer
Seq.\ \textbf{15} (2012), Article 12.3.8.
\bibitem{MO320838} MathOverflow, \emph{Can the product of two sums of prime reciprocals be $1$?},
question 320838 (comments by R.~Israel and J.~Rosen, answer by ``Woett''), 2019.
\url{https://mathoverflow.net/questions/320838}.
\bibitem{ErdosProblems307} T.~F.~Bloom, \emph{Erd\H{o}s problem \#307},
\url{https://www.erdosproblems.com/307}.
\bibitem{OEIS_A003415} OEIS Foundation, \emph{Sequence A003415 (arithmetic derivative)},
\url{https://oeis.org/A003415}.
\bibitem{teRiele1984} H.~J.~J.~te~Riele, \emph{Computation of all the amicable pairs below
$10^{10}$}, Math.\ Comp.\ \textbf{47} (1986), 361--368.
\bibitem{BHP2001} R.~C.~Baker, G.~Harman, and J.~Pintz, \emph{The difference between consecutive
primes, II}, Proc.\ London Math.\ Soc.\ (3) \textbf{83} (2001), 532--562.
\bibitem{Iwaniec1978} H.~Iwaniec, \emph{Almost-primes represented by quadratic polynomials},
Invent.\ Math.\ \textbf{47} (1978), 171--188.
\bibitem{LemkeOliver2012} R.~J.~Lemke Oliver, \emph{Almost-primes represented by quadratic
polynomials}, Acta Arith.\ \textbf{151} (2012), 241--261.
\bibitem{HardyWright} G.~H.~Hardy and E.~M.~Wright, \emph{An Introduction to the Theory of Numbers},
6th ed., Oxford Univ.\ Press, 2008.
\bibitem{Tenenbaum} G.~Tenenbaum, \emph{Introduction to Analytic and Probabilistic Number Theory},
3rd ed., Grad.\ Stud.\ Math.\ \textbf{163}, Amer.\ Math.\ Soc., 2015 (Erd\H{o}s--Wintner theorem).
\bibitem{Stothers} W.~W.~Stothers, \emph{Polynomial identities and Hauptmoduln}, Quart.\ J.\ Math.\
Oxford (2) \textbf{32} (1981), 349--370; R.~C.~Mason, \emph{Diophantine Equations over Function
Fields}, Cambridge Univ.\ Press, 1984.
\bibitem{Seva323194} Seva (\texttt{mathoverflow.net} user), \emph{A recursion with a number--theoretic
function}, MathOverflow question 323194 (2019). \url{https://mathoverflow.net/questions/323194}.
\bibitem{Bado2026} I.~O.~Bado, \emph{Structural constraints for an Erd\H{o}s unit--fraction problem
over primes}, preprint, Aix--Marseille University (2026).
\bibitem{Giuga1950} G.~Giuga, \emph{Su una presumibile propriet\`a caratteristica dei numeri primi},
Ist.\ Lombardo Sci.\ Lett.\ Rend.\ Cl.\ Sci.\ Mat.\ Nat.\ (3) \textbf{14} (1950), 511--528.
\bibitem{BBBG1996} D.~Borwein, J.~M.~Borwein, P.~B.~Borwein, and R.~Girgensohn, \emph{Giuga's conjecture
on primality}, Amer.\ Math.\ Monthly \textbf{103} (1996), no.~1, 40--50.
\bibitem{BorweinWong1997} J.~M.~Borwein and E.~Wong, \emph{A survey of results relating to Giuga's
conjecture on primality}, CRM Proc.\ Lecture Notes \textbf{11} (1997), 13--27.
\bibitem{BJM2000} W.~Butske, L.~M.~Jaje, and D.~R.~Mayernik, \emph{On the equation
$\sum_{p\mid N}1/p+1/N=1$, pseudoperfect numbers, and perfectly weighted graphs}, Math.\ Comp.\
\textbf{69} (2000), no.~229, 407--420.
\bibitem{GrauOllerMarcen2012} J.~M.~Grau and A.~M.~Oller-Marc\'en, \emph{Giuga numbers and the
arithmetic derivative}, J.\ Integer Seq.\ \textbf{15} (2012), Article 12.4.1; arXiv:1103.2298.
\bibitem{SondowMacMillan2017} J.~Sondow and K.~MacMillan, \emph{Primary pseudoperfect numbers,
arithmetic progressions, and the Erd\H{o}s--Moser equation}, Amer.\ Math.\ Monthly \textbf{124} (2017),
no.~3, 232--240; arXiv:1812.06566.
\bibitem{GrauOllerMarcenSadornil2021} J.~M.~Grau, A.~M.~Oller-Marc\'en, and D.~Sadornil,
\emph{On $\mu$-Sondow numbers}, arXiv:2111.14211 (2021); Acta Math.\ Hungar.\ (2022).
\bibitem{Cox} D.~A.~Cox, \emph{Primes of the Form $x^2+ny^2$: Fermat, Class Field Theory, and Complex
Multiplication}, 2nd ed., Wiley, 2013.
\bibitem{OEIS_A054377} OEIS Foundation, \emph{Sequence A054377 (primary pseudoperfect numbers)},
\url{https://oeis.org/A054377}.
\bibitem{OEIS_A007850} OEIS Foundation, \emph{Sequence A007850 (Giuga numbers)},
\url{https://oeis.org/A007850}.
\bibitem{Evertse1984} J.-H.~Evertse, \emph{On sums of $S$-units and linear recurrences}, Compositio
Math.\ \textbf{53} (1984), 225--244.
\bibitem{ESS2002} J.-H.~Evertse, H.~P.~Schlickewei, and W.~M.~Schmidt, \emph{Linear equations in
variables which lie in a multiplicative group}, Ann.\ of Math.\ (2) \textbf{155} (2002), 807--836.
\end{thebibliography}

\end{document}
